\documentclass[aps,amsmath,twocolumn,amssymb,floatfix,showpacs,superscriptaddress,nofootinbib,longbibliography]{revtex4-1}
\usepackage{MnSymbol}
\usepackage{braket}
\usepackage[dvipsnames]{xcolor}
\usepackage{float}
\usepackage{subfigure}
\usepackage{tikz}
\usepackage{comment}
\usepackage{textcomp}
\usepackage{physics}
\usepackage{lipsum}
\usepackage[colorlinks=true,linktoc=page,linkcolor=blue,citecolor=blue,urlcolor=blue]{hyperref}
%\usepackage[colorlinks=true,linktoc=page,linkcolor=blue,citecolor=blue,urlcolor=blue]
\newcommand*\circled[1]{\tikz[baseline=(char.base)]{\node[shape=circle,draw,inner sep=1.0pt] (char) {#1};}} % Circled character
\newcommand{\vect}[1]{\boldsymbol{\mathrm{#1}}}
\newcommand{\Eq}[1]{Eq.~(\ref{#1})}
\mathchardef\mhyphen="2D % Define a "math hyphen"
\newcommand{\noin}{\noindent}
\newcommand{\etal}{{\it et al.~}}
\newcommand{\ie}{{ i.e.,\,\,}}
\newcommand{\eg}{{e.g.,~}}
\newcommand{\ua}{{\uparrow }}
\newcommand{\da}{{\downarrow }}
\newcommand{\eps}{{\epsilon}}
\newcommand{\la}{{\langle}}
\newcommand{\ra}{{\rangle}}
\newcommand{\suc}{{superconductor~}}
\newcommand\bea{\begin{eqnarray}}
	\newcommand\eea{\end{eqnarray}}
\newcommand\beq{\begin{equation}}  
	\newcommand\eeq{\end{equation}}
\newcommand{\new}{\newpage}
\newcommand{\noi}{\noindent}
\newcommand{\bib}{\bibitem}
\newcommand{\cosec}{\operatorname{cosec}}
\newcommand{\non}{\nonumber}  
\usepackage[normalem]{ulem}
\definecolor{lime}{HTML}{A6CE39}
\usepackage{sidecap,tikz}
\DeclareRobustCommand{\orcidicon}{\hspace{-1.0mm}
	\begin{tikzpicture}
		\draw[lime, fill=lime] (0.0,0.0) 
		circle [radius=0.15] 
		node[white] {{\fontfamily{qag}\selectfont \tiny \,ID}};
		\draw[white, fill=white] (-0.0525,0.095) 
		circle [radius=0.007];
	\end{tikzpicture}
	\hspace{-3.0mm}
}
\foreach \x in {A, ..., Z}{\expandafter\xdef\csname orcid\x\endcsname{\noexpand\href{https://orcid.org/\csname orcidauthor\x\endcsname}{\noexpand\orcidicon}}
}

\newcommand{\orcidauthorA}{0000-0000-0000-0000}
\newcommand{\orcidauthorB}{0000-0001-8261-0483}

\AtBeginDocument{
	\newwrite\bibnotes
	\def\bibnotesext{Notes.bib}
	\immediate\openout\bibnotes=\jobname\bibnotesext
	\immediate\write\bibnotes{@CONTROL{REVTEX41Control}}
	\immediate\write\bibnotes{@CONTROL{%
			apsrev41Control,author="08",editor="1",pages="1",title="1",year="1"}}
	\if@filesw
	\immediate\write\@auxout{\string\citation{apsrev41Control}}%
	\fi
}

\begin{document}
%\title{Investigation of Kitaev-Heisenberg-$\Gamma$ model in small clusters}  
\title{Deciphering  competing  interactions of Kitaev-Heisenberg-\texorpdfstring{$\Gamma$}~ system
in  clusters}	

\author{Sheikh Moonsun Pervez\orcidA{}}
\email{moonsun@iopb.res.in}
\affiliation{Institute of Physics, Sachivalaya Marg, Bhubaneswar-751005, India}
\affiliation{Homi Bhabha National Institute, Training School Complex, Anushakti Nagar, Mumbai 400094, India}
%\author{Nicola Seriani\orcidB{}}
%\email{nseriani@ictp.it}
%\affiliation{The Abdus Salam ICTP, Strada Costiera 11, I-34151 Trieste, Italy}
\author{Saptarshi Mandal\orcidB{}}
\email{saptarshi@iopb.res.in}
\affiliation{Institute of Physics, Sachivalaya Marg, Bhubaneswar-751005, India}
\affiliation{Homi Bhabha National Institute, Training School Complex, Anushakti Nagar, Mumbai 400094, India}

%************************************** ABSTRACT ************************************%
\begin{abstract}
We study Kitaev-Heisenberg-$\Gamma$ Hamiltonian on finite size clusters and investigate  various zero and finite temperature properties in detail and explore the key differences due to relative sign and magnitude of Kitaev($K$), Heisenberg $(J)$) and $\Gamma$ interaction in the presence of magnetic field. Nearest-neighbor spin-spin correlation function shows step like character when $J$ is positive only irrespective of sign of $K$. Similarly the low energy spectrum under magnetic field depends on the sign of $K$. Dynamics of Majorana fermions depend greatly on external magnetic field which can even stabilize the gauge fields to Kitaev limit in the presence of $J$. Majorana dynamics also shows that $K-J-\Gamma$ model to be a testbed to manipulate quantum speed limit.   For small magnetic field, magnetization shows a dome like structure for an intermediate temperature range for positive $K$. However a positive $J$ can cause this for negative $K$ also. The zero temperature magnetisation and susceptibility shows large dependence on external magnetic field and relative sign and magnitude of $K$ and $J$. Either a negative $K$ or $J$ is shown to give qualitatively similar susceptibility for different magnetic field only.  For the specific heat a negative $K$ or $J$ favors mostly single peak structure while a two peak structure is favored for positive $K$ or $J$, though the magnitude of magnetic field can causes re-entrant behavior also. We briefly discuss effect of $\Gamma$ and compare with recent experiments. 
 
%Here we study Kitaev-Heisenberg-$\Gamma$ Hamiltonian on finite size clusters and investigate  various zero and finite temperature properties in detail. Our analysis shows some key differences between ferromagnetic and antiferromagnetic
%Kitaev($K$)(and Heisenberg $(J)$) interaction in the presence of magnetic field. For example nearest-neighbor spin-spin correlation function shows step like character when $J$ is positive only irrespective of sign of $K$ . The low energy spectrum behaves similarly across all clusters under magnetic field for negative $K$ but depends largely on cluster sizes for positive $K$. Our study shows that dynamics of Majorana fermions depend greatly on external magnetic field and depending on its magnitude and direction gauge fields can be stabilized to Kitaev limit. For small magnetic field, there exist an intermediate temperature regime where magnetization shows dome like structure for positive $K$ only when $J$ is zero. However a positive $J$ can make this happen for negative $K$ also. The zero temperature magnetisation and susceptibility shows large dependence on external magnetic field and relative sign and magnitude of $K$ and $J$.  Either a negative $K$ or $J$ is shown to give qualitatively similar susceptibility for different magnetic field only.  For the specific heat it is observed that either a negative $K$ or $J$ favors mostly single peak structure while a two peak structure is favored  for positive $K$ or $J$, though the magnitude of magnetic field can causes re-entrant behavior also.  
 

\end{abstract}
%************************************** ABSTRACT ************************************%
\date{\today}

%************************************** TITLE ************************************%	
\maketitle
%************************************** TITLE ************************************%



%************************************** INTRODUCTION ************************************%
\section{Introduction}
 Starting from its inception Kitaev model\cite{kitaev-2006} has witnessed a remarkable journey which is still continuing with overwhelming excitement. Though initially it was proposed as a mere toy model with certain aspect of quantum computation in perspective, it has now traversed almost all the aspect a quantum many body system can offer. To mention a few of them, exact short range  spin-spin correlation~\cite{smandal-2007} established its spin liquid nature, effect of disorder and localization~\cite{kao-2021,nasu-2020}has been investigated for such spin liquid. Realization of Kondo effects and its competition with Kitaev physics has been investigated~\cite{tathagata-2020} and spin-spin correlation function for perturbed  Kitaev model ~\cite{tikhonov-2011,lunkin-2019} has found a power low decaying spin-spin correlation function. The dynamics of perturbed Kitaev model have explored several aspects such as confinement-deconfinement transition,  quantum phase transitions in perturbed Kitaev model~\cite{saptarshi-2011,animesh-2020,animesh-2021, knolle-2018}. Kitaev model for classical 
spins~\cite{baskaran-2008,samarth-2010} and in variety of lattices in two and three dimensions \cite{kivelson-2007,vala-csl-2010,naveen-2008,eschmann-2020} showed interesting feature such as order-disorder mechanism, existence of chiral spin liquid, gapless fermi contour etc. \\
\indent
Apart from the above theoretically motivated works, akin to a dream realized in a mysterious way, very soon a plethora of materials have been proposed to posses Kitaev like interaction along with other kind of interactions ~\cite{chalaupka-2010,simon-2022,takagi-2019,motome-2020,hermanns-2018,abanerjee-2016}. The salient feature of these materials are that  the presence of non-Kitaev interactions make the coveted spin-liquid state non-existent at very low temperature. However at higher temperature the Kitaev physics mainly the fractionalization of spins into Majorana fermions and fluxes and the intrinsic differences in the length and time scale of their dynamics enable the spin liquid state to be revived at intermediate temperature\cite{janssen-2017,chern-2021,czajka-2021,niravkumar-2019,wulferding-2020,berke-prb-2020,
haoxiang-nature-2021,balz-prb-2021,plessis-2020}. Further recent studies show that proximity effect of $\alpha-RuCl_3$ in graphene heterostructure~\cite{sananda-2019,sananda-2021} increases the strength of Kitaev interaction and moreover induces an insulator to metal transition which might lead to exotic superconducting states. Such heterostructure study  holds a great promise for interesting Kitaev related physics
to be observed in near future. \\
\indent

One of the fascinating aspect of condensed matter system is that physical interactions at microscopic ranges govern the physical properties at macroscopic length \cite{anderson,fifty-years}. Commonly all interactions are extended over few atoms or molecules governed by mainly electrostatic interactions combined with quantum mechanical laws and inherent symmetries. It is also observed how the physical properties of the system changes as the size of the system varies from microscopic to macroscopic scale. One such example is quantum spin system.  Various phenomena such as long range order, symmetry breaking, fractionalization of spins, anyonic excitations, certain properties of spin-spin correlation functions, topological properties are 
only realized over large enough system sizes\cite{savary-review}. However in this regards, there are some exception in Kitaev Model where one of the defining property of Kitaev system i.e bond-dependent short range correlation function and spin-fractionalization is realized exactly even for the smallest size defined properly.  This renders some aspect of Kitaev magnetism or exactly for Kitaev model in small clusters \cite{tathagata-2020,kataoka-jpsj-2020}. Further the primary aspect of frustration or absence of it is mainly governed by local considerations of exchange coupling of spin-spin interaction though globally they drive the system completely differently. However it is an interesting question to itself how these local dynamics of commonly prevalent local exchange interaction competes with Kitaev type exchange interactions and whether there is a definite pattern in it~\cite{wang-2020}. To the best our knowledge such detail exact comparative study is absent within a single work. This motivates us to study Kitaev model in small clusters to explore several aspects in detail. Further due to inherent difficulties to deal with an interacting many body system it is not possible to to analysis in a single study various competing mechanism exist in Kitaev systems. For example the relative sign and strength of Kitaev interaction and Heisenberg interaction and $\Gamma$ interaction would make the analysis very demanding for a single study. To this end we consider exact analysis of Kitaev-Heisenberg-$\Gamma$ systems in small clusters and investigate some very  essential quantities such as eigenspectrum, nearest-neighbor spin-spin correlation function, flux operators, detail dynamics of gauge and matter Majorana fermions, specific heat, susceptibility and specific heat for all possible combinations of Kitaev interaction, Heisenberg interaction, $\Gamma$ and external magnetic field in $\hat{z}$ direction. We hope that our study would serve
as a valuable reference for Kitaev-Heisenberg-$\Gamma$ system. \\
%It is might have shed interesting light to understand Ktaev physics in general as well as to understand some  Kitaev spin liquid material $\rm{OsxCl_3}$ comprised of Honeycomb Nano-Domains \cite{kataoka-jpsj-2020}. \\
\indent
\begin{figure}[H]
\centering
\includegraphics[width=0.85\columnwidth,height=!]{schematics.pdf}
\caption{Schematics of Kitaev clusters with (a) 4-site, (b) 6-site, (c) 8-site, and, (d) 12-site are shown in panel a, b, c and d respectively. The red bonds, blue bonds and green dashed bonds  are having $x-x$-type, $y-y$  type and $z-z$ type Kitaev interaction respectively.}
\label{clusterfig}
\end{figure}
Our plan of presentation is as follows. We begin by investigating how the spectrum of Kitaev-Heisenberg-$\Gamma$ system evolves as we vary the strength of Heisenberg interaction  for various values of $\Gamma$ for different cluster of sizes in Sec. \ref{spectrum}. We also considered the effect of external magnetic field on the low energy eigenvalues. Our analysis shows that the eigenenergies of the Kitaev-Heisenberg-$\Gamma$ system van be grouped into two branches. In the first branch,  energies vary linearly for all ranges of Heisenberg coupling $J$ and for the second branch energies vary linearly only for large magnitude of $J$ but for small values they shows a cusp or dip near $J=0$ whose exact position depends on the value of $\Gamma$ and cluster sizes. Interestingly the linear branches changes their energies from positive to negative when the sign of $J$ is reversed but for second type it remains negative or positive as we change the sign of $J$. The incidental crossing between the two kind of energies can be controlled by tuning $\Gamma$. \\
\indent 
\begin{figure*}
	\includegraphics[width=1\columnwidth,height=!]{eigenvalues.pdf}
	\caption{Eigenvalues are plotted for various size of clusters as shown. The eigenvalues are renormalized by dividing with $6 \times N$  where $N$ is the number of  sites in a given cluster. First, second, third and fourth columns are for four site, six site, eight site and twelve site clusters as denoted.  The first three rows are for antiferromagnetic Kitaev interaction with $K=1$ and the last three rows are for ferromagnetic Kitaev interaction in Kitaev clusters for varying $J$ with fixed $K=-1$. The values of $\Gamma$ taken as shown inside each panel.}
\label{eigenspectrum}
\end{figure*}
   
After investigating the spectrum, we focus on the nearest-neighbor correlation function in Sec. \ref{correlation} and study how it depends on cluster sizes and on different system parameters as before. We have considered both $\sigma^z_i \sigma^z_j$ and $\sigma^y_i \sigma^y_j$. We find that correlation function shows interesting variations
depending on the relative sign and magnitude of Kitaev interaction ($K$), Heisenberg interaction $J$ and $\Gamma$.
Very interestingly we find a clear distinction between the ferromagnetic and antiferromagnetic Heisenberg interaction on the nearest neighbor correlation function. For the antiferromagnetic Heisenberg interaction the
correlation function shows plateau like behavior where spin-spin correlation does not change as we change $J$.
This suggest a competing mechanism between the quantum entanglement induced by Kitaev interaction and antiferromagnetic Heisenberg interaction. On the other hand for ferromagnetic Heisenberg interaction such plateaus are absent and shows various non-monotonous behavior showing a different competition between the quantum fluctuation by Kitaev interaction in response to ferromagnetic Heisenberg interaction which tries to induce a collinear order in the system. \\
\indent
  

Next we study the fate of conserved plaquette operator in Sec. \ref{bpoperator} and dynamics of Majorana fermions by exact numerical computation in Sec. \ref{majodynamic}. We find that the amplitude and time period of oscillations of Majorana fermions greatly depend on the relative sign of the Kitaev interaction ($K$) and Heisenberg interaction $(J)$. On an average a ferromagnetic Kitaev interaction is seen to increase the time period of oscillations. We noticed that the external magnetic field can modulate the nature of the dynamics of the Majorana fermion greatly and depending on the directions of it, the external magnetic field can stabilize the fluctuation of conserved gauge fermions. In certain cases it shows when the probability of Majorana fermion is zero describing a scenario to define quantum speed limit of Majorana fermions in these system. \\                                                                                                                                                             
\indent



Finally we present finite temperature magnetization, susceptibility and specific heat in Sec. \ref{magnetization}, \ref{susceptibility}. \ref{specificheat} respectively. We give a detail description how these quantities depend on
 various magnetic exchange coupling ($K, J, \Gamma$) including external magnetic field $h_z$.  The primary findings
on magnetization is that for antiferromagnetic Kitaev interaction and low magnetic field, there is an intermediate
range of temperature where magnetization posses a smooth  dome like behavior which is absent for ferromagnetic Kitaev interaction. However depending on the Heisenberg interaction this dome like structure can be re-entrant for ferromagnetic Kitaev interaction. The external magnetic field is shown to control the zero temperature magnetization which also depends differently for different combinations of $J, K$ and $\Gamma$. Similar to magnetization, we find that ferromagnetic and antiferromagnetic Kitaev interaction has distinct features which again can change differently
depending on the sign and magnitude of Heisenberg interaction and applied external magnetic field. \\
\indent

\begin{figure*}
	\includegraphics[width=1\columnwidth,height=!]{eigenvalues_with_hz_at_Gamma0.pdf}
	\caption{Few lowest eigenvalues are plotted with  external magnetic field applied in $z$-direction. Each column denotes a particular cluster as mentioned at the top. Each row denotes a combination of Kitaev ($K$) coupling and Heisenberg coupling (J) as mentioned at the right side of each row. The ground state energy ($E_0$) is shown by green dots and first excited state $E_1$ is shown by blue bots and few more excited eigenstates $E_n$ with $n>1$ are shown by red dots.}
\label{eigenvalue-in-hz}
\end{figure*}
\section{Kitaev-cluster-Basics and spectrum analysis}
\label{spectrum}
The defining Hamiltonian we investigate in our study is $J-K-\Gamma$ model\cite{wang-2021} which
is given as,
\begin{eqnarray}
&& H= \sum_{\langle i,j \rangle} \left(J {\bf S}_i \cdot {\bf S}_j + K S^{\gamma}_iS^{\gamma}_j + \Gamma(S^{\alpha}_iS^{\beta}_j + S^{\beta}_iS^{\alpha}_j)  \right)
\end{eqnarray}



We begin our study of various sizes of Kitaev cluster by first analyzing the eigenvalue spectrum and understand qualitatively the effect of Heisenberg and $\Gamma$ interaction as plotted in Fig. \ref{eigenspectrum}. We start with four site cluster and discuss $\Gamma=0$ at first. We notice that eigenvalues could be grouped into two categories. Firstly there are eigenstates whose energies are linearly varying with $J$ and secondly there are eigenstates whose energies vary non-linearly  for small values of $J$ but eventually vary linearly for large values of $J$. This is true for all clusters. We find that a finite $\Gamma$ affects the eigenspectra in two ways. Firstly it breaks the degenerates of eigenvalues found for $\Gamma=0$. It is interesting to note that only those eigenvalues which vary linearly become non-degenerate but the resultant eigenenergies again vary linearly with $J$. Secondly the band width (the separation between the highest and lowest eigenvalues) increases for non-zero $\Gamma$ and higher for positive $\Gamma$. Also it is evident that $\Gamma$ determines the gap at $J=0$. The physical reason why $\Gamma$ interaction breaks the degeneracy of the eigenvalues is that it breaks the reflection symmetry under $\sigma_x \rightarrow - \sigma_x$ but $\sigma_{y,z} \rightarrow \sigma_{y,z}$. \\
%\textcolor{red}{we  may plot the gap versus $\Gamma$ plot}. More interestingly the closes negative eigenvalue at $J=0$ may belong to eigenvalue with positive or negative slope.
\indent

For large positive $J$, the four site cluster has large singlet($\frac{1}{2}\ket{\uparrow_1\downarrow_2-\downarrow_1\uparrow_2}\times\ket{\uparrow_3\downarrow_4-\downarrow_3\uparrow_4}$) contribution to the ground state. The other states which are varying linearly can be approximated by classical state such as $\left( \uparrow_1\downarrow_2\uparrow_3\downarrow_4\right) $ having energy $E_{cl}=-2J+2K$ and $\left( \uparrow_1\downarrow_2\downarrow_3\uparrow_4\right) $ ($\left( \uparrow_1\uparrow_2\downarrow_3\downarrow_4\right)  $) having energy $E_{cl}=-2J-2K$. The highest (lowest) eigenvalues for positive (negative) $J$ can be approximated as $E_{cl}=6J+2K$ ($E_{cl}=-6J+2K$) which is determined  by an approximated ferromagnetic state though quantum fluctuation is more effective for antiferromagnetic $J$ to increase the magnitude of the eigenvalue  further. It is evident that as we decrease the value of $J$ from a very large positive value to a large negative value, the quantum ground state  is continually adjusted from a large singlet component to a large ferromagnetic component for $\Gamma=0$. However for non-zero $\Gamma$, as we go from large antiferromagnetic $J$ to large ferromagnetic $J$ there is a crossover between the states $E_{0}$ to $E^n$ at some finite $J$. Here $E_0$ is the lowest energy state at $J \rightarrow - \infty$ and $E^{n}$ is some positive high energy states at $J \rightarrow - \infty$. These two states continually come closer to each  other and for $\Gamma=0$, the asymptotically merge but for finite $\Gamma$ there is a crossover between these two states at some critical value of $J$.   \\
\indent
In the appendix, we enlist all the eigenvalues and eigenstates for four site cluster for finite $K, J$ and $\Gamma=0$ for future references. We also yield exact expression for correlation function for comparison in the appendix. To complete our discussion on eigenvalue spectra on finite size cluster  we note that the general features present in four site cluster prevail for larger site clusters also. Mainly there exists two group of eigenstates one with linearly varying eigenvalues and one
with non-linearly varying with $J$. We always find the number of linearly varying eigenvalues constitute the majority of eigenstates and thus determine the finite temperature property. \\
In previous studies a lot of attention as given to the effect of magnetic field on the Kitaev interaction in general and it is observed that it stabilizes the Kitaev spin-liquid phase at intermediate temperature \cite{berke-prb-2020,nasu-motome-2015,koga-nasu-2019,li-2021-ncom}. Hence it is important to  examine what intricacies an external magnetic field brings in to the eigenspectrum. We plotted the few low energy eigenstates for various clusters as shown in Fig. \ref{eigenvalue-in-hz}. \\
\begin{figure*}
	\includegraphics*[width=1.2\columnwidth,height=!]{zz_correlations.pdf}
	\caption{Nearest neighbor correlation on a $z$-type bond in Kitaev clusters in the presence of magnetic field $\vec{h}$ in the $\hat{z}$ direction is plotted.  A given column denotes a particular cluster as mentioned on the top. Each row is plotted for a particular $(K,J)$ as written at the right side. Only specific magnetic field with representative characteristics has been plotted though magnetic field has been taken from 0.2 to 2 insteps of 0.2. The color legends for magnetic field is shown at the left.}
\label{zzcorrelation}
\end{figure*}
For the sake of simplicity we first start analysing the effect of external magnetic field for $J=0$ and $J=-0.5$
as shown by the first, second and third and fourth rows respectively. We observe that for these two choices of
Heisenberg interaction, when Kitaev interaction is ferromagnetic the eigenvalues mainly increase with magnitude
of $h_z$ as evident from the second and fourth rows. However for the antiferromagnetic Kitaev interaction the eigenvalues decrease or increase depending on the cluster size and strength of external magnetic field. Notably for twelve site cluster for $(J,K)= (-0.5,-1)$ (panel `$(\tilde{l})$') we observe the gap between the ground state
and first excited states are gradually increasing with $h_z$ but for antiferromagnetic Kitaev interaction
i.e for $(J,K)= (-0.5,1)$ (panel `$(l)$') the ground state and the first excited states remain nearly degenerate.
Also in the former case the other excited states move away from the ground state energy but in the later the low energy excited states move closer toward the ground state energy. This highlights that the relative sign and magnitudes of Kitaev and Heisenberg interaction compete very differently. Similar analysis can also be done for antiferromagnetic Heisenberg interactions. It is interesting to compare our result to that found in an earlier study~\cite{zhu-2018} where effect of an external magnetic field in the direction of $(1,1,1)$ has been considered. Ferromagnetic and antiferromagnetic Kitaev interaction have been shown to produce qualitatively different results. In the case of ferromagnetic Kitaev interaction the first excited  state and other excited states shows monotonous increase of the gap from the ground state energy which we also observe as evident from first row Fig. \ref{eigenvalue-in-hz}. However for the antiferromagnetic Kitaev interaction a qualitatively similar evolution of eigenenergies have been observed where higher excited states come closer to the ground state energy and at some critical external field it becomes nearly degenerate with the ground state \cite{zhu-2018}. From the first row of Fig. \ref{eigenvalue-in-hz} we observe a qualitatively different evolution of low lying eigenstates with external magnetic field  than found for ferromagnetic Kitaev interactions. However the difference between the earlier result~\cite{zhu-2018} and in our case  is attributed to the finite size\cite{holzmann-2016} and the direction of magnetic field which has been taken in $(0,0,1)$ directions. 

\section{Correlation function}
\label{correlation}

In the pure Kitaev limit the two-spin correlation function vanishes beyond nearest neighbor and also they are bond dependent i.e on $\alpha$-bond only $\langle s^{\alpha}_i s^{\alpha}_j \rangle$ exists, a consequence of $Z_2$ conserved operators. On the other hand in the pure Heisenberg limit, the spin-spin correlation function for nearest neighbor spins are isotropic i.e all $\langle s^{\alpha}_i s^{\alpha}_j \rangle$ exist for any $\alpha=x,y,z$ and they are equal, a consequences of $SU2$ symmetry. The characteristics bond-dependent correlation function in the Kitaev limit do not survive when  any other interaction such as Heisenberg or $\Gamma$ or external magnetic field is applied. As the nature of the bond-dependent short-range correlation function at Kitaev limit does not depend on the system size and it is realized in  exact fashion even for small clusters, it is thus natural to ask what is the manifestation of other non-Kitaev interaction such as Heisenberg interaction, $\Gamma$ interaction and effect of external magnetic field. The previous theoretical and many experimental findings have shown presence of a plethora of phases for such generalized Kitaev-Heisenberg-$\Gamma$ model\cite{khwang-2022,feng-2022,han-li-2023,yilmaz-2022,shuyi-2023,smit-2020,kimchi-2014,yamaji-2016,young-baek-kim-2015,wang-2021}. There are  long range ordered various collinear phase, Kitaev-spin liquid phase and other quantum paramagnet phase. It is generally understood that long-ranged classically ordered collinear order arises mainly by local constraint to minimize the bond-interactions which has some reciprocity with the Kitaev interaction that here also the local $Z_2$ conserved operator induced constraint are the deciding factor for short range correlation function. Thus we think it is indeed of great relevance to investigate the correlation function even for small clusters. However it is known that we will not be able to infer the complete information regarding long-ranged correlation brought in by non-Kitaev interaction such that the spin-liquid phase has long-ranged algebraically decaying correlation function.  Nevertheless the correlation function study for the small cluster promises to show some signature of competition between different kind of interactions. To begin with we first discuss the scenario $\Gamma=0$.  \\ 
\indent
\subsection{$z-z$ correlation function}
Without loss of generality we consider the bonds $(1,3)$, $(1,4)$, $(1,8)$ $(1,12)$ for the four site, six site, eight and twelve site cluster. We note that they are all $z$-type bond as shown in Fig. \ref{clusterfig}. \\
\indent
{\it Correlation function for $\Gamma=0$:} Let us first analyze the correlation function for $\Gamma=0$ and $h_z=0$. As we see  from the red line plot that for ferromagnetic or antiferromagnetic  $J$, the $z-z$ correlation function monotonically  increases (decreases) for ferromagnetic(antiferromagnetic) Kitaev interaction. This causes a point of inflexion at $J=0$. However we observe that for four site cluster the correlation function changes drastically with external magnetic field and become independent of whether Kitaev interaction is ferromagnetic or antiferromagnetic. In both the cases the correlation function saturates to $\pm 1$ for antiferromagnetic or ferromagnetic Heisenberg interaction. However for $K=1$ as we increase the value of the external magnetic field we find that after a critical magnetic field $h_{z,c}$, the $z-z$ spin correlation function  remains constant for certain range of antiferromagnetic Heisenberg interaction. The ground state does not changes much  such that   spin-spin correlation function remains constant for a range of Heisenberg interaction.  The above nature for $z-z$ correlation function changes as we increase the cluster sizes with two notable differences. Firstly the zero field correlation function for antiferromagnetic Heisenberg interaction yields a saturated $z-z$ correlation function. For  ferromagnetic interaction the $z-z$ correlation function monotonically increases and reach a saturation whose upper bound is 1 for very large external magnetic field. However, interestingly, for very large external magnetic field, the spin-correlation-plateau is universally observed for all cluster for certain ranges of antiferromagnetic Heisenberg interaction for example consider $(N,K,h_z)=(6,1,2), (6,-1,2),(12,1,2), (12,-1,2)$. In short the main differences between the ferromagnetic and antiferromagnetic Heisenberg interaction on the nearest-neighbor $z-z$ correlation function is that for the former the $z-z$ correlation function behaves mostly monotonously but for the later the $z-z$ correlation function shows plateau like characteristics where $z-z$ correlation function remains constant for certain range of antiferromagnetic Heisenberg interaction and the width of the plateaus  increase with increasing magnetic field.\\
\indent
\begin{figure*}
	\includegraphics*[width=1.2\columnwidth,height=!]{yy_correlations.pdf}
	\caption{Nearest neighbor $y-y$ correlation on a $z$-type bond in Kitaev clusters in the presence of magnetic field $\vec{h}$ in the $\hat{z}$ direction is plotted.  A given column denotes a particular cluster as mentioned on the top. Each row is plotted for a particular $(K,J)$ as written at the right side. Only specific magnetic field with representative characteristics has been plotted though magnetic field has been taken from 0.2 to 2 insteps of 0.2. The color legends for magnetic field is shown at the left.}
\label{yycorrelation}
\end{figure*}
{\it Correlation function for $\Gamma=\pm 0.5 $:} For the four size cluster, we observe that for antiferromagnetic Kitaev interaction plateau-like  property continues to exist though for ferromagnetic Kitaev interaction it ceases to exist. Otherwise the $z-z$ spin correlation function mostly increases (decreases)
as we increase the external magnetic field for antiferromagnetic (ferromagnetic) Heisenberg interactions. For six and eight  site cluster the behavior of $z-z$ spin correlation function are qualitatively similar to that of $\Gamma=0$, the plateau like behavior  persists for multiple ranges of antiferromagnetic Kitaev interactions and occasional oscillations appear near $J=0$. Apart from this occasional oscillations which may arise due to degenerate ground states, the $\Gamma$ interaction brings in two important differences. First we note that for $\Gamma=0$, the asymptotic value of $z-z$ correlation function for large Heisenberg  interaction is independent of the strength of  external magnetic field but in the presence of $\Gamma$ the asymptotic value of $z-z$ correlation function for ferromagnetic Heisenberg interaction depends on the external magnetic field and increases with the strength of it. Secondly we find, for $K=1$, as we increase the cluster size, the $\Gamma$ interaction brings the $z-z$ correlation function to lower value for some small range of ferromagnetic Heisenberg interaction such that for twelve site cluster the $z-z$ correlation function decreases for small values of  $h_z$. This initial decrease is more pronounced in the presence of $\Gamma$ interaction. Thus we see multiple effects being in competition such as quantum correlation due to increased size of the cluster as well as due to the competition between Heisenberg and $\Gamma$ interaction.  We note that for $\Gamma=-0.5$, the $z-z$ spin-spin correlation function is qualitatively similar to that of $\Gamma=0.5$, though for certain values of Heisenberg interaction there is some differences. For example for  small values of $h_z$ and for ferromagnetic Heisenberg interaction we observe a value of $J_c$ for which the $z-z$ correlation function shows a maxima, this is pronounced for six site cluster and eight and twelve site cluster for some intermediate value of ferromagnetic $J$.
\begin{figure*}
	\includegraphics[width=1.2\columnwidth,height=8cm,keepaspectratio]{Bp.pdf}
	\caption{ Various plaquette operators are plotted against Heisenberg interaction $J$. Each column denotes a particular cluster as mentioned at the top. First row and second row correspond to antiferromagnetic and ferromagnetic Kitaev interaction with $K= 1$ and $K=-1$ respectively as mentioned at the right side. The color legends show the value of magnetic field taken. Inside each panel of the first row, the expression for particular plaquette operator is given.}
\label{bpfigure}
\end{figure*}
%\begin{figure*}
%	\includegraphics*[width=1.2\columnwidth,height=!]{papfig/yy_correlations.pdf}
%	\caption{\textcolor{red}{caption will be updated after latest figure available}Nearest neighbour correlation on a $z$-type bond in Kitaev clusters in the presence of magnetic field $\vect{h}$ in the $\hat{z}$ direction.  Zero temperature spin-spin correlations are measured for varying $J:-3\rightarrow+3$. $z$-$z$ correlations with $K=+1.0$ are shown for (a$_N$) $\Gamma=-0.5$, (b$_N$) $\Gamma=0.0$, (c$_N$) $\Gamma=+0.5$ and $y$-$y$ correlations with $K=+1.0$ are shown for (d$_N$) $\Gamma=-0.5$, (e$_N$) $\Gamma=0.0$, (f$_N$) $\Gamma=+0.5$ in $N$-site cluster, where $N$ = 4, 6, 8, $\&$ 12. The right three columns correspond to $K=-1.0$ and the nomenclature is same but with a $\tilde{~}$ sign (example, \~{a}$_N$).}
%\end{figure*}
\subsection{$y-y$ correlation function} In the Kitaev limit $y-y$ correlation function on a $z$-type bond has zero value. However the Heisenberg and $\Gamma$ interaction brings in some non-zero values to $y-y$ correlation functions. When $\Gamma=0, h_z=0$, $y-y$ correlation function
saturates to some constant positive (negative) value for very large ferromagnetic or antiferromagnetic Heisenberg interaction. For intermediate ferromagnetic Heisenberg interaction there are some non-monotonicity due to the competition of Kitaev and Heisenberg interaction. As we apply external field $h_z$, the asymptotic value of $y-y$ correlation function is zero for ferromagnetic Heisenberg interaction indicating that ground state is polarized in the $z$-direction. But the asymptotic $y-y$ correlation for the large antiferromagnetic Heisenberg interaction is finite even for large external magnetic field. This suggests that the ground state is still a quantum ground state and this qualitatively marks the difference between the effect of ferromagnetic and antiferromagnetic Heisenberg interaction to the Kitaev interaction.\\
\indent
For intermediate strength of Heisenberg interaction interesting variations of $y-y$ correlation appear which is counter intuitive. For ferromagnetic Heisenberg interaction  the $y-y$ correlation function first increase (opposite to the naive expectation) with $h_z$ being non-zero which is remarkable. Because naively we do not expect $y-y$ correlation to gain in strength due to a field applied in $z$-direction. This suggests the quantum nature of the ground state for intermediate strength of ferromagnetic Heisenberg interaction  and as it is departing from Kitaev limit the Heisenberg interaction is manifesting itself in non-zero large $y-y$ correlation function. As we increase the strength of ferromagnetic Heisenberg interaction, the $y-y$ correlation function starts to decrease.  This results in a critical $J_{c,fm}$ for each $h_z$ where $y-y$ correlation function is maximum. For the antiferromagnetic Heisenberg interaction we observe the plateau like behavior for small values of $J$. The width of the plateaus increase as we increase the strength of external magnetic field. It is important to note that this plateau like behavior persists even for twelve site cluster.\\

\indent
   Study of correlation function has attracted a lot of previous interest \cite{tikhonov-2011,song-2016,lunkin-2019,yang-2020,berke-prb-2020} where the spin-spin correlation function has been calculated in the presence of an external magnetic field or other interactions. However either these studies examined the long time and long distance behavior of correlation function \cite{tikhonov-2011,lunkin-2019} or approximately for short distances \cite{yang-2020} perturbatively. On the other hand our exact analysis  uncovered interesting nature of correlations due to competition of Heisenberg and Kitaev interaction. This reveals that the qualitative nature of nearest neighbor correlation function for  antiferromagnetic  Heisenberg interaction  is completely different than ferromagnetic Heisenberg interaction in the presence of external magnetic field which has not been explored before.
%\begin{figure*}
%	\includegraphics[width=1.2\columnwidth,height=!]{papfig/Bp.pdf}
%	\caption{$B_p$ for varying Heisenberg iteraction $J:-3.0\rightarrow+3.0$ with fixed magnetic field $h_z$ with $K=+1$ in (a) 4site with triangular plaquette, (b) 4site with four-sided plaquette, (c) 6-site, (d) 8-site, and (e) 12-site cluster. $B_p$ for $K=-1$ are shown in (f) 4site with triangular plaquette, (g) 4site with four-sided plaquette, (h) 6-site, (i) 8-site, and (j) 12-site cluster.}
%\end{figure*}
\section{Expectation of plaquette conserved operator}
\label{bpoperator}

The exact solvability of Kitaev model resides on the fact that it reduces to non-interacting Majorana fermion hopping problem due to presence of conserved plaquette operators which manifest itself as local $Z_2$ gauge field on a given bond. These $Z_2$ gauge fields are bilinear in   Majorana fermions. We note that these plaquette conserved operators are product of multispin operators and cover almost same size of the cluster for the four and six site clusters. In the Kitaev limit they are conserved and static but in the presence of non-Kitaev interactions such as Heisenberg and external magnetic field they are no longer conserved and  acquire dynamics. The $B_p$ operators can always be expressed as $B_p= \prod \sigma^{\alpha_1}_{1}... \sigma^{\alpha_n}_{n}$ where `$1$' to `$n$' denote the indices of sites on the plaquette in clockwise or anticlockwise manner and $\alpha_{m}$ denotes the $m$-type outgoing Kitaev bond at the site $m$. For the Kitaev cluster defined on a tetrahedra there are three sites in an elementary plaquette while there are a four site conserved operator as well which is given by $B_{1-4}=\sigma^z_1 \sigma^z_2 \sigma^z_3 \sigma^z_4$. For the six and eight site cluster also the elementary plaquettes are squares. On the other hand for the twelve site cluster the elementary plaquettes are hexagons. The eigenvalue of the plaquette conserved operators are $\pm 1$ as evident from the identity $B^2_p=1$. In the Kitaev limit when $ J=0=h_z$, the eigenstates of the Hamiltonian are simultaneous eigenstates of $B_p$'s. However for the ground state, the eigenvalue of $B_p$ could be $+1$ or $-1$ depending on the total number of sites in the cluster.  Various symmetries such as time reversal symmetry, reflection symmetry also determine the ground state eigenvalues of $B_p$ for various clusters~\cite{lieb-1994}.   \\
\indent
%\begin{figure*}
%	\includegraphics[width=2\columnwidth,height=!]{project1_Bp.pdf}
%	\caption{$B_p$ for varying Heisenberg iteraction $J:-3.0\rightarrow+3.0$ with fixed magnetic field $h_z$ with $K=+1$ in (a) 4site with triangular plaquette, (b) 4site with four-sided plaquette, (c) 6-site, (d) 8-site, and (e) 12-site cluster. $B_p$ for $K=-1$ are shown in (f) 4site with triangular plaquette, (g) 4site with four-sided plaquette, (h) 6-site, (i) 8-site, and (j) 12-site cluster.}
%\end{figure*}
\begin{figure*}
\includegraphics[width=2.4\columnwidth,height=10cm,keepaspectratio]{gauge-majorana-dynamics.pdf}
\caption{Dynamic correlation function $|g(t)|$ is plotted as a function of time $t$ for $\Gamma=0$.
Each column represents Majorana dynamics for various cluster as mentioned at the top. In the right we denoted
the Heisenberg strength $J$ for which $|g(t)|$ has been plotted. Odd rows and even rows represent plots for antiferromagnetic and ferromagnetic Kitaev interaction respectively. Various color legends explain the magnitude of external magnetic field $h_z$.}
\label{gauge-majo-dyn}
\end{figure*}
 We begin our discussion with tetrahedral system as shown in Fig. \ref{bpfigure}, panel $(a)$ and $(f)$ for $K=1$ and $K=-1$ respectively. For this system ground state is degenerate with $B_p=\pm 1$  as being product of three spin operator it is odd under time reversal symmetry.  However we choose $B_p=1$ for conveniences which has been used in Fig. \ref{bpfigure} for $J=0, h_z=0$. However for non-zero $J$ such a state selection is not possible and the exact ground state eigenfunction determines  its expectation value. The value of $B_p$ being zero at finite $J$ for $K=\pm 1$ denotes that ground state eigenfunction is equal superposition of $B_p=\pm 1$. This can also be explained on the basis that Hamiltonian is time reversal invariant but the $B_p$ operators are not and thus its expectation value is zero. In panel $(b)$ and $(g)$ of Fig. \ref{bpfigure} we plotted expectation of $B_{1-4}$
which can be obtained by the multiplication of two plaquette operators associated with triangles. Thus it can be thought of as a correlation between the two adjacent plaquette operators associated with two triangles. We observe that at $J=0$, its value is $+1$ by our choice. But for antiferromagnetic Heisenberg interaction it changes its value to -1 implying that the two $B_p$ associated with triangles are anti-correlated and this constant value increases with increasing magnetic field. After a critical $J_{afm}$, $S_{1-4}$ reduces to zero implying that the conserved quantities associated with triangles are no longer correlated. It manifests an unique competition between the quantum paramagnet states pertaining to Kitaev limit and quantum paramagnet states in the Heisenberg limit.\\
\begin{comment} \textcolor{blue}{need to revisit the plot as it is contradictory to Lieb theorem, what is the reason for it.The four site cluster with square plaquettes shows $B_p=-1$ due to Lieb theorem owing to refelction symmetry. It is interesting that for finite $h_z$ it stabilizes the $B_p$
to -1 (because it Zeeman term $h_z\sigma^z_i$ commutes with $B_p$) even for non zero values of $J$. Why for $K=-1$ $B_p=0$ at $J=0$ and -1 at finite J only, needs to discuss. However this peculiar behaviour of changing $B_p$ from 1 to -1 can be used for quantum gate operation.} 
\end{comment}
\indent
For the six site clusters the  $B_p$ shows $+1$ for $J=0, h_z=0$. There are six four site plaquettes and one six site plaquette and here we have plotted the $B_p$ associated with the four site plaquette. Each two four site plaquette yields the six-site plaquette. Due to absence of reflection symmetry for the elementary square we find the expectation value of plaqeutte operator as 1 instead of -1 (which should be according to Lieb theorem~\cite{lieb-1994} if reflection symmetry is present). However for  the six site plaquette operator we obtain the expectation value $1$ at the Kitaev limit i.e $h_z=0=J$.  For nonzero $J$ and $K=1$ we observe that there is a window of $J$ for which  $B_p$ expectation takes a constant positive value and it goes to zero after a critical ferromagnetic $J$ whose magnitude increases with external magnetic field. For ferromagnetic Kitaev interaction this critical Heisenberg interaction shift to antiferromagnetic $J$  suggesting an interesting competition between the relative sign of $J$ and $K$.\\
\indent
Having described four site and six site clusters in detail we now explain the ground state $B_p$ expectation value for eight site and twelve site cluster. We note that the elementary plaquettes in eight site cluster is four and hence Lieb theorem gives $B_p=-1$ for the ground state. We find that for $J=0, h_z=0$ $B_p$ to be $-1$ as expected as seen from Fig. \ref{bpfigure}. For $K=1$ any large values of $\pm J$,  $B_p$ vanishes as expected because of eigenstates mixing with all eigenstates of $B_p$ in the presence of $J$. However for small value
of $\pm J$ and $K=1$, there are high values of $B_p$ for some values of $h_z$. This probably indicate different plaquette operator transform differently under external magnetic field. We note that in the eight site cluster there are plaquettes operators which contains only $\sigma^z$ and hence commutes with the external magnetic field. Thus there are preferences for some eigenstates to be selected and this causes finite large $B_p$ for small $\pm J$. However there is a remarkable differences for ferromagnetic Kitaev interaction in comparison to antiferromagnetic Kitav interaction. For $K=-1$ for large $-J$ the $B_p$ saturates for $h_z=0$ and also the half-width maxima of $B_p$ for other $h_z$ is broader in comparison to antiferromagnetic $K$. For the 12 site cluster we observe that for $J=0, h_z=0$ and $K=\pm 1$  $B_p=1$ satisfying Lieb's theorem. For finite $h_z$ and antiferromagnetic $K$, $B_p$ shows a rapid oscillations (having positive and negative values of $B_p$) for small window of $|J|$ and vanishes thereafter.  The same is true for ferromagnetic $K$ though for small value of $J_z$, the $B_p$  shows a saturation for large ferromagnetic $J$ \cite{zhu-2018,han-li-2023,nandini-2019,liang-2018,yilmaz-2022}.\\
\indent
\begin{comment}
\textcolor{red}{ Having discribed the qualitative feature of variations of $B_p$ with Hisenberg interaction we now try to explain recent experimental and other theoretical observations. To this end it may be important to investigate the effect of changing the direction of $B_p$ and find whether there is oscillations of $B_p$ with this and systematic observation how does it vary for ferromagnetic and antiferromagnetic $J$. This might help to understand the gauge field dynamics and contribution to spin-liquid at finite temperature and oscillations with the angle dependence of exernal magnetic field}
\end{comment}

\section{Dynamics of Majorana fermions}
\label{majodynamic}
Much of the recent interest in explaining various experimental observation and theoretical results was devoted to understand the gauge field dynamics in Kitaev-Heisenberg-$\Gamma$ model. It is well known that a pure Kitaev model is described by a non-interacting Majorana fermion hopping problem ($H= \sum_{ij} J_{ij} i c_i c_j $) in the background of conserved $Z_2$ gauge fields $u_{ij}= ib_i b_j$. In appropriate representation these conserved  gauge field operators can be re-written as an equivalent number operator as $u_{ij}= (2 \chi^{\dagger} \chi -1 )$, named bond-fermion \cite{mandaljpa}. Also the Majorana operator ($c_i$)  can also be appropriately recasted interms of complex fermion $\psi$ which on some bonds  appear as $2 \psi^{\dagger}_i \psi_i-1$. Depending on the interest there were much effort  investigating the gauge field and matter field dynamics by studying the expectation value $<i c(t)_i c(0)_j>$ and also $<i b_i(t)b_j>$ in the presence of other non-Kitaev interaction such as Heisenberg interaction, $\Gamma$ interaction and external magnetic field \cite{knolle-2018,nasu-motome-2019}. These external perturbations make the gauge fields to acquire a dynamic as they are no longer conserved and a quantitative analysis of these matter field and gauge field dynamics enable us to understand the possible spin-fractionalization leading to spin-liquid phase appearing at intermediate temperature \cite{berke-prb-2020,nasu-motome-2015,koga-nasu-2019,li-2021-ncom}. This is also intimately connected to confinement-deconfinement transitions associated with magnetically ordered state to quantum paramagnetic states including Kitaev spin liquid state. To study the gauge field and matter field dynamics here we use  much simpler quantities ($g_{\phi}(t)~ {\rm and}~ g_{\chi}(t) $) defined below which are equally capable  to characterize  the dynamics of Majorana fermions brought in  by non-Kitaev interaction. 

\begin{eqnarray}
\label{gphipsi}
&&g_{\phi}(t)= \langle N_{\phi,i}(t) N_{\phi,i}(0) \rangle, ~~~ N_{\phi,i}= \phi^{\dagger}_i \phi_i \nonumber \\
&&g_{\chi}(t)= \langle N_{\chi,i}(t) N_{\chi,i}(0) \rangle, ~~~ N_{\chi,i}= \chi^{\dagger}_i \chi_i
\end{eqnarray}

The reason that we choose this operator is that $N_{\phi,i}$ ($N_{\chi,i}$)acts as a projector on the ground state such that it projects only states with `$\phi$' ($\chi$) fermion on the bond `$i$'. Thus $g_{\phi}(t) $ (or $g_{\chi}(t)$) yields the probability amplitude of those states to return back to itself. In the Kitaev limit as $N_{\chi,i}$ is static and thus  $g_{\chi}(t)$ is constant in time but acquire time dependence once non-Kitaev terms are considered. In the presence of non-Kitaev  interactions $\psi$ and $\chi$ interact with each other thus estimating those quantities in Eq.~\ref{gphipsi} yield an understanding of the dynamics associated with non-Kitaev interactions. In Fig. \ref{gauge-majo-dyn} and Fig. \ref{matter-majo-dyn}, we show how  $g_{\chi}(t)$ and $g_{\phi}(t)$ depends on time for various cluster and other system parameters for $\Gamma=0$. Below we summarize some salient characteristics of gauge and matter Majorana fermion dynamics. \\

%\begin{figure*}
%\label{majo-dyn}
%\includegraphics[width=1.2\columnwidth,height=!]{papfig/gauge-majorana-dynamics.pdf}
%\caption{$|g(t)|$ is examined as a function of time $t$ with $K=1.0,~\Gamma=0.0,~h_z=0.0$ and for $J=0.0, \& 0.5$ in (a) 4-site, (b) 6-site, (c) 8-site, and in (d) 12-site cluster. In pure Kitaev limit, the gauge fermion remains static over time; when $J\neq0$, it becomes dynamic. The matter fermion is already interacting with each other even in the absense of Heisenberg interaction, and $|g(t)|$ shows dynamical behaviour for them, as expected.}
%\end{figure*}
\begin{figure*}
\includegraphics[width=1.2\columnwidth,height=8cm, keepaspectratio]{matter-majorana-dynamics.pdf}
\caption{Dynamics of matter Majorana fermion as defnied by the quantity $g_{\phi}(t)$ is plotted for various cluster size and parameter values. Various panels describe a particular cluster, a given combination of Kitaev interaction strength $K$, Heisenberg interaction strength $J$ as described in Fig. \ref{gauge-majo-dyn}.}
\label{matter-majo-dyn}
\end{figure*}
\subsection{Dynamics of matter and gauge fermions at $\Gamma=0$}
To describe the dynamics of majorana fermions we consider bond (1,3), (1,4), (1,4), (1,6) for four site, six site, eight site and twelve site cluster. Thus in the below text whenever we mention $\chi$ or $\phi$ fermion, it is to be understood that they are defined on this bonds. Occasionally for eight site and twelve site cluster we mention $\chi_2$ and $\phi_2$
and they are defined on bonds (2,7) and (2,9) for eight and twelve site cluster respectively.
{\small {\it {Qualitative differences between four site cluster and other clusters:}}}
In all the parameter regime there are qualitative differences in the dynamics of gauge fermions between the four site cluster and other larger clusters. For $\Gamma=0, J=0$ the dynamics of gauge fermions are very regular, periodic as we vary the external magnetic field. For intermediate antiferromagnetic $J$, the oscillations are very regular and amplitude decreases monotonously as we increase the strength of $h_z$ (see Fig. \ref{gauge-majo-dyn}, panel `$i$'),  though for ferromagnetic $J$ appearances of additional frequency could be noticed which points out that the ground state is modified differently than the antiferromagnetic case.  \\
\indent
The same qualitative difference as discussed above can be seen for matter fermions also. However we notice that in general the average amplitude of oscillations for gauge fermions are always larger than the matter fermions. This can be attributed to non-local nature of the gauge fermions interms of the spin-operators and also to the fact that each flux sectors are separated by an energy gap which is much larger than the fermionic spectrum for each flux sector. \\
\indent
{\small {\it {Regular oscillations and irregular oscillations:}}}
It is observed that depending on the strength of external magnetic field the oscillation of gauge fermions could be regular to irregular type. For example consider six site cluster $J=0$. Up to $h_z=0.6$, the oscillations are very small but after that it suddenly increases\cite{yamada-2020} and reaches maximum for $h_z=1.2$ and then again it decreases, this has been shown in Fig. \ref{gauge-majo-dyn}, panel $(b)$. In numerous occasions it have been observed that $h_z$ can reverse the response of gauge fermions amplitude or time period  of oscillations.  On the other hand for eight site cluster a different scenario is observed comparing
$J=-0.5, 0, 0.5$ i.e Fig. \ref{gauge-majo-dyn}, panel $(k)$. For $J=-0.5$ the amplitude of oscillation decreases with increase of $h_z$, this behavior is opposite to that found for $J=0$ and for $J=0.5$, the oscillations remain more or less same for all $h_z$ (see Fig. \ref{gauge-majo-dyn}, panel $(g)$ ), no special feature is seen.  However for twelve site cluster such non-monotonous response is absent. very interestingly we see from Fig. \ref{gauge-majo-dyn} panel $(c), (d), \tilde{c}, \tilde{d}$ that the gauge fermion does not show any dependency on time for any magnetic field. This can be referred as stabilization of gauge fields under external magnetic field on certain bonds. However on a different bonds the gauge fermion dynamics is oscillatory as shown by extra four panel at the right side of the main panel. The above controlled manipulation of gauge field dynamics could be of practical use.\\
\indent
The identical trends can be seen for matter fermions as well where depending on the relative strength of external magnetic field $h_z$ and Heisenberg interaction $J$. For example one can compare $(N,K,J,h_z)=(6,-1,0,0.4)~ {\rm and}~ (6,-1,0,1.2)$. In the former case there are many minimum found but for the later only a single minima found, see Fig. \ref{matter-majo-dyn}, panel $(b)$, the nature of oscillation completely changed. Similarly for $(N,K,J,h_z)=(6,-1,0.5,0), (6,-1,0.5,0.4), (6,-1,0.5,0.6)$ panel $(\tilde{f})$ in Fig. \ref{matter-majo-dyn}(similar trends are seen  for $\phi_2$ fermion as well). These all three set of parameters yield three different behavior for $\phi_2$ fermion. For eight site cluster such differences are also found for $(K,J)=(1,-0.5), (1,0), (1,0.5)$ (see Fig. \ref{matter-majo-dyn}, panel $(\tilde{c}),(\tilde{g}),(\tilde{k})$ and also for $(K,J)=(-1,0), (-1,0.5)$ where depending on external magnetic field $\phi_1$ fermions dynamic changes a lot (see Fig. \ref{matter-majo-dyn}, panel $(k)$). Similarly for twelve site clusters the relevant set of parameters for such change in the nature of $\phi_2$ fermions can be seen for $(K,J)=(1,-0.5), (1,0)$. \\
\indent


{\small {\it {Dependency of average amplitude and time period of oscillations on $h_z$ and $J$:}}}
Interestingly it has been noticed that for $J=0$, the time period of oscillations is much larger  for $K=-1$ than $K=1$ case, this can be noticed from the first and second rows of Fig. \ref{matter-majo-dyn}. In the first row, for finite $h_z$, the oscillations has more amplitude and time period than the second row.  However as we turn on finite $J$ the time period of oscillations become comparable though it greatly depends on the relative strength of $K, J$ and $h_z$.  It is also observed that the external magnetic field helps in reviving the amplitude of matter fermion on that site. For example consider twelve site cluster where the rapid fluctuations are absent in comparison to other smaller cluster. If we look at $(K,J)=(-1,0.5)$, Fig. \ref{matter-majo-dyn}, panel $(\tilde{h})$  where in between $t=1$ and $t=3$, there is a revival of amplitude of matter fermions. Similarly we find for   $ (K, J)=(-1,0), (\pm1,-0.5)$ in panel $(d)$ with $h_z=0.4$ and panel $(\tilde{l})$ and $(l)$. Similarly we see that for gauge fermions, $K=-1, J=0$ corresponds to larger time period in comparison to $K=1, J=0$ (see Fig. \ref{gauge-majo-dyn} first and second row ).\\

\indent
{\small {\it {Nature of oscillations for large $\pm |J|$:}}}
Nature of oscillations in large $J$ (ferromagnetic of antiferromagnetic) limit is more regular. Yet depending on the relative sign of $J$ and $K$ there could be presence of one single frequency or more. This is reflected in the amplitude also. In some cases amplitude is constant when there is a single frequency and amplitude varies when multiple frequencies are present. For four site cluster large positive $J$ the oscillation  contains single frequency (independent of $h_z$) though large $-J$ contains multiple frequency centered around two main frequencies for non-zero $h_z$. For six site cluster the oscillations are still regular though multiple frequencies appear but for eight site cluster it becomes regular due to special symmetry of the cluster. For twelve site cluster there are presence of many frequencies but they vary slowly with $h_z$ so that the oscillations appear coherent. For matter fermion, the amplitude of oscillations are always smaller in comparison to gauge fermions except special cases where the external magnetic fields are seen to stabilize the gauge fields on some particular bond. \\
\indent

{\small {\it {Two different channels of dynamics of oscillations:}}}
It is also noted that depending on initial position of the  matter fermions its time dependency could be very different. For comparison we note the results for twelve site cluster with the parameter value $(K,J)= (-1, 0.5)$. We observe that whether $\phi_1$ is decaying with rapidly, $\phi_2$ is having quite small oscillations and vary smoothly. The average magnitude of $g_{\phi}(t)$ corresponding to $\phi_2$ also gradually increases consistently with increasing $h_z$. This different behavior depending on the initial position of the matter fermion indicates
the local competition between the different interaction such as Heisenberg and Kitaev. The same is true for ferromagnetic Heisenberg interaction with parameter $(K,J)=(-1,-0.5)$. There is a marked difference on the dynamics of $phi_1$ and $\phi_2$. The same can be said for $J=0$. This observation could be very much useful for the controlled manipulation of Majorana fermions for useful quantum state operations.\\
%\begin{figure*}
%\includegraphics[width=1.2\columnwidth,height=!]{papfig/matter-majorana-dynamics.pdf}
%\caption{$|g(t)|$ is examined as a function of time $t$ with $K=1.0,~\Gamma=0.0,~h_z=0.0$ and for $J=0.0, \& 0.5$ in (a) 4-site, (b) 6-site, (c) 8-site, and in (d) 12-site cluster. In pure Kitaev limit, the gauge fermion remains static over time; when $J\neq0$, it becomes dynamic. The matter fermion is already interacting with each other even in the absense of Heisenberg interaction, and $|g(t)|$ shows dynamical behaviour for them, as expected.}
%\label{matter-majo-dyn}
%\end{figure*}
\indent
{\small {\it {Stabilization of $Z_2$ gauge field by external magnetic field:}}} It has been observed that as the applied magnetic field is in $z$-direction it can stabilize the $Z_2$ gauge field\cite{khwang-2022} as seen in for eight site cluster  ($K,J$)= ($-1, -2$) or ($-1,-0.5$), Fig. \ref{gauge-majo-dyn}, panel $(\tilde{k})$ and also for $\tilde{l}$ and other). For intermediate value of $J$ positive  the gauge field fluctuation decreases as we increase $h_z$ as seen for ($K,J$)= ($-1, 0.5$), Fig. \ref{gauge-majo-dyn}, panel $(\tilde{f})$. This can be thought of as stabilizing the spin-liquid or fractionalized state by external magnetic field. This is also true for $K=1$ though there are some fluctuations of the gauge field around values closed to unity. Identical trends are seen for twelve site cluster, a finite $h_z$ greatly stabilizes $Z_2$ gauge fields on certain bonds. This can be seen for $(K,J)= (1,-2), (1,-0.5)$ (Fig. \ref{gauge-majo-dyn}, panel $(l)$) and $(K,J)= (-1,-2), (-1,-0.5)$(Fig. \ref{gauge-majo-dyn}, panel $(\tilde{l})$). For antiferromagnetic $J$, the $Z_2$ gauge fields show more oscillatory pattern in comparison to ferromagnetic $J$, they become closer to unity. \\
\indent
{\small {\it {Scenarios to define quantum speed limit:}}} Recently various aspect of quantum speed limit as defined by the evolution of a quantum state from a given state to a distinct different state attracted wide interest for various reasons \cite{deffner-2017,gal-ness-2022,pati-2022,arindam-2022}. In this context we find that quantum-speed limit could be investigated as far as the evolution of Majorana fermions are considered. Initially they begin with a given occupancy at a given bond and if at any future time, the occupation of the fermion on that bond becomes zero, we can define quantum speed limit for our system. In our study of dynamics of gauge and matter fermion we witnessed numerous occasions where such situations occur. For reference purpose
we denote by $(N,K,J,h_z)$ the set of parameter where $N$ denotes the number of sites in the cluster, $K$ and $J$ are the strength of Kitaev and Heisenberg interaction respectively and $h_z$ being the strength of external magnetic field applied in $z$-direction. \\
\indent
Most fascinating occurrence of quantum speed limit happens for four site cluster with $(N,K,J,h_z)=(4,-1,0,2), (4,-1,0.5,0.8)$. This happens in both the matter and gauge fermion sector. For the second set of parameters they are  remarkably repeating with time. For $K=-1, J=-0.5$ the  quantum speed limit phenomena can also be observed depending on the strength of external magnetic field. On the other hand for antiferromagnetic Kitaev interaction such realization does not happen except for very large ferromagnetic Heisenberg interaction. From this we note that when the sign of Kitaev interaction and Heisenberg interaction is opposite then only it happens. \\
For the six site cluster and for the gauge fermions we find $(N,K,J,h_z)= (6,1,0,1.2)$ being such an example. For large antiferromagnetic $J$ it also happens for matter fermions. For eight site cluster also there are occasions for $\phi_2$ with $(N,K,J,h_z)= (8,1,0,1.6), (8,1,0,1.4), (8,1,0,1.8)$. There are occasions for gauge fermions also for such phenomena as well as for ferromagnetic Kitaev interactions for example  $(N,K,J,h_z)= (8,-1,0,0.6)$. For twelve site cluster also there are special cases where such phenomena are observed. A careful study and estimation of this quantum speed limit will be done in future study as this is not the primary focus here. \\
\indent
\begin{figure*}
	\includegraphics[width=1.5\columnwidth,height=!]{magnetization_in_zero_field.pdf}
	\caption{ Zero field spontaneous magnetization has been plotted for various cluster. Each row denotes a particular cluster as denoted at the top. At the right side of each row we mention the value of $(K, \Gamma)$ taken. Magnetisation has been plotted for various values of Heisenberg interaction ($J$) as denoted by the color legends at the left.}
\label{magnetizationhzzero}
\end{figure*}
{\small {\it {The main differences in the dynamics of majorana fermions between twelve site cluster and other clusters:}}} The most obvious difference between the twelve site and other smaller cluster can be understood by analyzing $J=0$. As we see that the gauge fermion remains constant (for cluster of all sizes) with time but the matter fermion oscillates with a fixed frequency and returns back to unity after characteristic time for four site and eight site cluster due to the back flow from the boundary. On the other hand for six site cluster the matter fermions return back to unity after few local maxima though for twelve site cluster the matter fermions never returns back to unity. Secondly  for finite $J$, the oscillations in $\chi$ or $\phi$ fermions are having regular pattern for different  $h_z$ but for smaller cluster the pattern of dynamics vary greatly depending on the magnitude of the external magnetic field.\\


\indent
The appearance of multiple frequencies in the dynamics of Majorana fermions can be understood as follows. For a fixed gauge field configurations, the effective Hamiltonian for $\psi$ fermions reduce to a superconducting Hamiltonian with the presence of $ \eta_{ij} \phi_i \phi_j + h.c$. The ground state $|G \rangle$ is constructed as the vacuum of the quasi-particles where the states with even number of $\psi$ fermions (including zero) appear with different probability amplitude. Now the operator $N_{\phi,i}= \phi^{\dagger}_i \phi_i$ selects (or projets) only those states  with the dimer `$i$' being occupied. These daughter states could be re-expressed as linear combinations of all the eigenstates (including excited and ground states) with different amplitudes. Thus the final expression of  $g_{\phi}(t)$ contains the summation of $ \sum_{n} A_{n}\exp^{-(E_n-E_0)}$ where the detail analytic expression of $A_n$ is in general very complex. With system size very large, a large number of contribution to $A_{n}$ coming from numerous daughter states makes the probability amplitude oscillate with no fixed time period. \\
\indent
Now for finite non-Kitaev interaction being present the gauge fermions are no longer conserved and hence they also oscillate with time. For the four site cluster due to special symmetry the matter fermion and gauge fermion returns back to unity with identical frequency.  As the cluster size is increased, the oscillations of matter and gauge fermions become  irregular  due to contributions from increasing number of excited states for larger clusters. However few universal characteristic features are present. One notes that the on an average  the gauge fermion amplitude is always larger than the matter fermion and closer to unity. Also matter fermions complete many cycles of oscillations within a given time window in comparison to gauge fermions. This points out that gauge fermions can be still approximated as slow varying background conserved operators for the matter fermions  and hence may support survive of decondined phase. \\
%\indent
%\subsection{Effect of finite $\Gamma$ on dynamics of gauge and matter fermions}
%\textcolor{red}{ To do it or leave it? Now we proceed to explain what happens for finite magnetic field and effect of $\Gamma$ interactions.}

\section{Magnetization}
\label{magnetization}
Having discussed various zero temperature properties including the eigen-spectrum analysis, nearest-neighbor correlation function, expectation of flux operators and detailed dynamics  of Majorana  fermions we now focus on some finite temperature properties of various clusters. Till now we have found that various competing interactions such as Kitaev interaction ($K$), Heisenberg interactions($J$) and $\Gamma$ interactions have effected the above mentioned  zero temperature and ground state properties in distinct ways. Thus it is natural to investigate  some  finite temperature properties to see what are the distinct ways each of those different interaction effect. Here we investigate magnetization, susceptibility and specific heat as described below.  
\subsection{Spontaneous magnetization}
%The clusters are subjected to temperature and a magnetic field $\vect{h}=h\hat{z}$ with $h=1$. Magnetization is measured along $\hat{z}$ direction. At a particular temperature $T=1/\beta$, magnetization per unit site is given by,
%\begin{eqnarray}\label{finite_T_magnetization}
%M^{(\hat{z})}(\beta)=\frac{1}{N}~\frac{\sum_{j=1}^{N}\sum_{l=1}^{2^N}e^{-\beta (E_l-E_0)}\expval{\sigma_j^z}{\varphi_l}}{\sum_{l=1}^{2^N}e^{-\beta (E_l-E_0)}},
%\end{eqnarray}
%where $\ket{\varphi_l}$ is the $l^{\text{th}}$ eigenstate of the Hamiltonian with eigenvalue $E_l$; $E_0$ is the ground state energy.
%\begin{figure*}
%	\includegraphics[width=1.5\columnwidth,height=!]{papfig/magnetization_in_zero_field.pdf}
%	\caption{Measurement of magnetization $M^{(\hat{z})}(\beta)$ for  different temperature. In the presence of $K=+1.0$, $M^{(\hat{z})}(\beta)$ is obtained for different values of $J$ with ($a_N$) $\Gamma=-0.5$, ($b_N$) $\Gamma=0.0$, ($c_N$) $\Gamma=+0.5$ in $N$ site cluster, where $N=4,~6,~8,~\&~12$. For $K=-1.0$, $M^{(\hat{z})}(\beta)$ is shown in graphs following same nomenclature but with a \~{}  sign (for example $\tilde{a}_N$).}
%\end{figure*}
\begin{figure*}
	\includegraphics[width=1.2\columnwidth,height=!]{magnetization_Gamma0_trimmed.pdf}
	\caption{Finite temperature magnetization has been plotted for various external magnetic field.  Each column describes result for  a specific cluster as mentioned at the top. Odd rows are for antiferromagnetic Kitaev interaction $K=1$ and the even rows are for ferromagnetic Kitaev interaction. The value of Heisenberg interaction taken for each rows are mentioned at the right of each row.}
\label{magnetizationfinitehz}
\end{figure*}

To calculate the spontaneous magnetization in the absence of external magnetic field, we considered a very small external magnetic field in the Hamiltonian and calculated the spontaneous magnetization along $\hat{z}$ direction. In the Fig. \ref{magnetizationhzzero} we have plotted the magnetization-vs-temperature for various parameters. In the below we describe some of the salient features.\\ 
\indent 

{\it \small{For ferromagnetic Kitaev interaction, $K=1$ :}} The zero field spontaneous magnetization reveals that for $\Gamma=0,~J=0$ the magnetization linearly increases as we decrease the temperature from a very high value. However there is an intermediate  temperature $T^*$ below which different scenarios may happen depending on the sign and magnitude of Heisenberg coupling $J$ as well as $\Gamma$ term. For $\Gamma=0$ and ferromagnetic $J$, the linear increase of magnetization with decreasing temperature saturates to a constant value at some critical temperature that we call $T_{c,fm}$. On the other hand for anti-ferromagnetic Heisenberg interaction, there exist two critical temperatures what we define as $T_{>,afm}$ and $T_{<,afm}$  with $ T_{<,afm} < T_{>,afm}$ such that below $T_{>,afm}$ the zero field magnetization starts to decrease  until $T_{<,afm}$ is reached and then it saturates to a constant value $\mathcal{M}$. We note that the constant magnetization $\mathcal{M}$ generally depends on the strength of magnitude of Heisenberg interaction and in general increase or decrease with the increase of magnitude of ferromagnetic or antiferromagnetic Heisenberg interaction. \\
Having explained the behavior of magnetization for $\Gamma=0$ we now describe what happens for finite $\Gamma$. We note that generally at high very high temperature, the magnetization is independent on the sign and magnitude of $\Gamma$. However below the critical temperature $T^*$ the effect of $\Gamma$ can be seen.
For ferromagnetic $\Gamma$ and $J$, the magnetization increases linearly with decrease in temperature and deviates from the saturating trends below $T^*$ for $\Gamma=0$. However for antiferromagnetic $\Gamma$ and ferromagnetic $J$ the the trends to reach  saturating magnetization remains intact though the value of saturation magnetization decrease with antiferromagnetic $\Gamma$. \\
Very interestingly for  antiferromagnetic $J$, the behaviors of magnetization seems to be indifferent of the sign of $\Gamma$. The dome like behavior of magnetization in between $T_{<,afm}$ and $T_{>,afm}$ exist though the saturation magnetization decreases in the presence of $\Gamma$ irrespective of its sign.\\
\indent
{\it \small{For antiferromagnetic Kitaev interaction, $K=1$ :}} The overall qualitative nature of variation of magnetization for antiferromagnetic Kitaev interaction in the presence of $J$ and $\Gamma$ with their possible signs and magnitudes are similar to those of ferromagnetic Kitaev interaction. All the critical temperature $T_{c,fm}, T_{<,afm}, T_{>,fm}$  are present characterizing the particular behavior of magnetization for anti-ferromagnetic Kitaev interaction. However it is found that for ferromagnetic $K$, the value of $T_{c}$ decreases as evident from comparing  Fig. \ref{magnetizationhzzero} ($b_8$) and Fig. \ref{magnetizationhzzero} ($\tilde{b}_8$) for eight site cluster. Similar observations are made for twelve site cluster by comparing  Fig. \ref{magnetizationhzzero} ($b_{12}$) and Fig. \ref{magnetizationhzzero} ($\tilde{b}_{12}$).  Though $T_{>}$ seems to be remain identical to that of ferromagnetic Kitaev interaction, the value of $T_{<}$ seems to decrease a little for ferromagnetic Kitaev interaction. The most remarkable difference between the ferromagnetic and antiferromagnetic Kitaev interaction is that the spread of saturation magnetization (or the band width of magnetization) at $T=0$ decrease for ferromagnetic interaction signifying the role of quantum fluctuation playing a dominant role for antiferromagnetic Kitaev interaction. Also we have noticed a decrease in bandwidth of magnetization at zero temperature as we increase the size of the cluster which is expected as due to quantum confinement quantum fluctuation plays a dominant role in smaller size clusters.   \\
%\textcolor{red}{\it \small {What can we say about fractionalization of spins at intermediate temperature }}

\begin{figure*}
	\includegraphics[width=1.2\columnwidth,height=!]{susceptibility_Gamma0_trimmed.pdf}
	\caption{Finite temperature susceptibility has been plotted for various external magnetic field.  Each column describes result for  a specific cluster as mentioned at the top. Odd rows are for antiferromagnetic Kitaev interaction $K=1$ and the even rows are for ferromagnetic Kitaev interaction. The value of Heisenberg interaction taken for each rows are mentioned at the right of each row. In the inset we have plotted susceptibility for $J=\pm 2$ as mentioned at the right of the figure.}
\label{susceptibilitygamma0}
\end{figure*}

\subsection{Magnetization in the presence of external field}

Having described the spontaneous magnetization at zero field, we now proceed to briefly describe the magnetization at finite external magnetic field. The finite field magnetization has been shown in Fig. \ref{magnetizationfinitehz}. Here we have taken external magnetic field ranging from $0.2$ to $2.0$ insteps of
$0.2$ but only those magnetization with a given external field has been shown for which there is a characteristic differences than the other. Below we describe some essential features. \\
\indent
{\it \small{Ferromagnetic ($K=-1$) vs antiferromagnetic Kitaev ($K=1$) for $J=0$:}} In the absence of Heisenberg interaction we find that there is a characteristic difference between the ferromagnetic and antiferromagnetic Kitaev interaction as evident from the first two rows of Fig. \ref{magnetizationfinitehz}. The dome like structure
appearing in the intermediate temperature in the absence of external magnetic field only appear for antiferromagnetic Kitaev interactions for small to intermediate magnetic fields. However for ferromagnetic Kitaev interaction the dome like behavior in the magnetization in the intermediate temperature are completely absent
even for small  external magnetic field. The other important difference is that the saturation magnetization for  given external magnetic field is higher for ferromagnetic Kitaev interaction than the antiferromagnetic Kitaev interaction. The saturation magnetization is almost independent of external magnetic field for ferromagnetic Kitaev interaction for all cluster. But for antiferromagnetic Kitaev interaction the saturation magnetization greatly depends on the cluster size and gradually increases with cluster size. \\
\indent
{\it \small{Ferromagnetic ($K=-1$) vs antiferromagnetic Kitaev ($K=1$) for $J=0.5$:}} In the presence of antiferromagnetic Kitaev interaction we find that when $K=1$ then the dome like structure in the magnetization is more pronounced and even exist for very high magnetic field. The most remarkable things for $K=1$
is that there are two limits of saturation magnetization $\mathcal{M}_1$ and $\mathcal{M}_2$ where $\mathcal{M}_2 > \mathcal{M}_1 \ge 0$. The low and intermediate magnetic fields yield a constant saturation magnetization $\mathcal{M}_1$ but the high magnetic field saturation magnetization yields another constant magnetization $\mathcal{M}_2$. On the other hand for ferromagnetic Kitaev interaction the saturation magnetization is different for each magnetic field though there is a re-appearances of dome like structure for small magnetic field. Thus
we see remarkable characteristic differences in the magnetization profile for ferromagnetic and antiferromagnetic Kitaev interaction and the fact that dome like structure in the magnetization reappear in the presence of antiferromagnetic Heisenberg interaction points out that the quantum frustration brought in by the antiferromagnetic Heisenberg interaction is responsible for that.\\
\indent
{\it \small{Ferromagnetic ($K=-1$) vs antiferromagnetic Kitaev ($K=1$) for $J=0.5$:}} For the ferromagnetic
Heisenberg interaction we observe that the zero temperature magnetization limit for $K=1$ depends on the external magnetic field and almost proportional to the external magnetic field. The dome like structure exist though becomes weaker. For the ferromagnetic kitaev interaction the zero temperature saturation magnetization is very high and becomes independent of external magnetic field after a critical external magnetic field. The dome like
structure for the intermediate temperature range is completely absent. Thus we see that effect of competing interactions manifested distinctly in the magnetization profile.





\section{Susceptibility}
\label{susceptibility}

Having described the behavior of magnetization in great detail we now proceed to discuss the susceptibility. As before we find characteristic signatures of competition among different interactions. In Fig. \ref{susceptibilitygamma0} we have plotted susceptibilities for various clusters and different combinations of parameters. As before though we have investigated susceptibility for finite magnetic field $h_z$ from $0.2$ to $2.0$, we have only plotted the representative one to describe how susceptibilities vary with magnetic fields. We fist begin with a comparison of susceptibility for $J=0$. \\
\indent
{\it \small{Ferromagnetic ($K=-1$) vs antiferromagnetic Kitaev ($K=1$) for $J=0$:}}
Let us first describe how the susceptibility varies when Heisenberg interaction is absent and what is the differences between ferromagnetic and antiferromagnetic Kitaev interactions. As we see from the second row
of the Fig. \ref{susceptibilitygamma0}, the susceptibilities for different clusters for ferromagnetic Kitaev
interaction are similar.  The susceptibility is zero at very high temperature and maximum at some intermediate 
temperature and becomes constant as temperature tends to zero. The peak height gradually decreases with
increasing magnetic field and also appears at higher temperature as expected. The saturation magnetization also depend on magnetic field and decreases with increasing magnetic field.\\
\indent
 However for antiferromagnetic Kitaev interaction the effect of external magnetic field on susceptibilities are completely different. First of all, there is no divergences in the susceptibilities  for any magnetic field which has been observed for ferromagnetic Kitaev interaction. Also the behavior of susceptibility is not uniform
for different magnetic field over the cluster sizes. For example for the twelve site cluster, the susceptibilities can be high or low at a given temperature depending on the external magnetic field. Identical trends can bee seen for eight site cluster. However though there are peaks in the susceptibilities at a given magnetic field for twelve site cluster, for eight site clusters peaks are very mild and only appear for small magnetic field. Owing to the small sizes and intricate nature of ground state, the susceptibility for four and six sites  cluster shows complex behavior.\\
\indent
{\it \small{Ferromagnetic ($K=-1$) vs antiferromagnetic Kitaev ($K=1$) for $J=0.5$:}} An antiferromagnetic
Heisenberg interaction brings in a very important difference to the susceptibilities of ferromagnetic Kitaev interaction which is that the  divergences found for $J=0$ at low magnetic field now appear at some finite magnetic field for twelve site cluster. On the other hand four and eight site cluster there are no pronounced maxima at any given magnetic field but such pronounced maxima appears for six site cluster. In a similar way we find that though in the absence of Heisenberg interaction there is no pronounced maxima for antiferromagnetic Kitaev interaction, an antiferromagnetic Heisenberg interaction causes susceptibility to have pronounced maxima or diverge at some large magnetic field. Also at low temperature the susceptibilities converge to very low values for antiferromagnetic Kitaev interaction but for ferromagnetic Kitaev interaction it saturates to quite large values. Thus we see clear evidences of competing effect on the susceptibilities due to Heisenberg and Kitaev interaction.  \\
\indent
{\it \small{Ferromagnetic ($K=-1$) vs antiferromagnetic Kitaev ($K=1$) for $J=-0.5$:}} In the presence of
ferromagnetic Heisenberg interaction we observe that for $K=-1$, the peak structure re-appears in susceptibilities and it is divergent mostly for small magnetic field. As we increase the magnetic field, the peak height gradually decreases across all clusters and at very low temperature they tend to saturate at zero for intermediate to large magnetic field as evident from the last raw of Fig. \ref{susceptibilitygamma0}. For the antiferromagnetic Kitaev interaction the pronounced peak appears for twelve site cluster at some intermediate critical magnetic field. Except this critical magnetic field, the susceptibilities are  generally regular though at a given temperature the  value of susceptibilities is not monotonously increasing or decreasing on external magnetic field. For four and six site cluster pronounced peaks re-appear at very small magnetic field but for eight site cluster  there is no peak like structure rather there is a mild dome like structures. Thus we observe that external magnetic field greatly influences the susceptibilities. These detail observations could be helpful in understanding the experiments in real material.
\begin{figure}
	\includegraphics[width=0.5\columnwidth,height=!]{sp-cartoon.pdf}
	\caption{A cartoon picture of the electronic spectrum in the Kitaev limit. Here $\rm{B_1, B_2, B_3, B_4}$ denote the flux configurations and the rectangular bar denote the eigenspectrum for a given flux configuration. The minimum of energy for each flux configurations as well as the separation  between two successive energy levels depends  on the flux configuration.}
\label{spcartoon}
\end{figure}
\begin{figure*}
	\includegraphics[width=1.2\columnwidth,height=!]{zero_field_Cv_plots.pdf}
	\caption{Normalized specific heat $C_v$  is plotted in the absence of magnetic field for various values of
Heisenberg interaction $(J)$ as shown inside the panel (a). Each column denotes a particular cluster as mentioned at the top. Solid lines denote antiferromagnetic Kitaev interaction and dashed lines denote the ferromagnetic Kitaev interaction with $K=1, -1$ respectively. }
\label{zerofieldcv}
\end{figure*}
\section{Specific heat}
\label{specificheat}
%Average Energy is given by,
%\begin{eqnarray}
%E(\beta)=\frac{\sum_{l=1}^{2^N}e^{-\beta (E_l-E_0)}E_l}{\sum_{l=1}^{2^N}e^{-\beta (E_l-E_0)}},
%\end{eqnarray}
%and, at temperature $T=1/\beta$, specific heat per unit site is given by,
%\begin{eqnarray}
%C_v=C(\beta)=-\frac{1}{N}~\beta^2~\frac{\partial E(\beta)}{\partial \beta}.
%\end{eqnarray}
We now proceed to analysis the specific heat for the Kitaev-Heisenberg-$\Gamma$ system. Much of the recent experimental results invested a great deal of analysis of some specific characteristics of specific heat
for this system as it can effectively signify the existence of spin-liquid phase or magnetically ordered states.
A theoretical scenario for eigenspectrum at the Kitaev limit is shown in Fig. \ref{spcartoon}. The effective Hamiltonian is described by various conserved flux configurations and each flux configurations yield a eigenspectrum as shown in Fig. \ref{spcartoon}. Generally the lowest energy for each flux configuration 
is different for different flux configurations. In the Fig. \ref{spcartoon}, we have denoted by $T_1$ and $T_2$,
two temperature such that $T_2 > T_1$. At low temperature as we see that the details of low energy eigenstates
contributions matter and hence the ground state energy of each  flux configuration dominate the contribution to specific heat yielding a first peak. However at very high temperature the low energy details of the spectrum for the each flux configurations are dominated by the high energy part of the eigenspectrum yielding a second peak where the gap between the minimum energies of each flux configuration does not matter.  This two gap structure of the specific heat drew a lot of attention as the deciding criteria for the spin-liquid phase or deconfined phase to be there. As we will observe that depending on the relative strength of Kitaev interaction $K$, Heisenberg interaction $J$, $\Gamma$ and external magnetic field $h_z$ we find existence of well established two peak structure, a single sharp peak or a wide broad peak. Primarily we assign the existence of clearly established flux sectors and the eigenspectrum above it when there are two well established peaks. On the other hand a single sharp peak would imply a strongly interacting confining phase with no emergent flux structure to be present. However the broad wide peak would either mean that two peak structure survive but they are close enough and hence can not be separated or the excitations are dominated by interacting magnon as found in recent studies.\\
%We examine specific heat of the clusters for ferromagnetic (FIG.~\ref{specific_heat_ferromagnetic_J}) Heisenberg interaction $J\leq0$ and for anti-ferromagnetic one (FIG.~\ref{specific_heat_antiferromagnetic_J}), $J\geq0$, with different $\Gamma$ values in the background of $K=1.0$. Reference:\cite{Kataoka_osxcl3,PhysRevLett.113.197205}.
%\begin{figure*}
%	\includegraphics[width=1.2\columnwidth,height=!]{papfig/zero_field_Cv_plots.pdf}
%	\caption{$C_v$ normalized by the cluster size. In the presence of $K=+1.0$, $C_v$ is obtained for different values of $J$ with ($a_N$) $\Gamma=-0.5$, ($b_N$) $\Gamma=0.0$, ($c_N$) $\Gamma=+0.5$ in $N$ site cluster, where $N=4,~6,~8,~\&~12$. For $K=-1.0$, $C_v$ is shown in graphs following same nomenclature but with a \~{}  sign (for example $\tilde{a}_N$). In all cases, specific heat tends to zero at $T\rightarrow\infty$ and at $T\rightarrow 0$, as expected.}
%\label{zerofieldcv}
%\end{figure*}
\indent
Now we discuss how the specific heat behaves for various values of $J$ and $\Gamma$ at zero magnetic field as shown in Fig. \ref{zerofieldcv}. 
%The specific  heat is an important measure of effective degrees of freedom  being activated at particular temperature and contributing to the thermodynamic quantities. From the essential physics of Kitaev interaction we know that at Kitaev limit fermionic excitations can be labeled by a a set of flux quantum number. The minimum energy for every flux combinations are usually different. Thus we notice that there are two energy scale in the system. First the ground state energies of each flux sector may form a certain energy scale in the prblem. On the otherhand the fermionic excitations over each flux combinations sets anothe energy scale.  The first energy scale we label as the flux sectors excitations which are basically composed of the low lying energy states over each flux sectors. And the second energy scale we label as fermionic excitations such that the intermediate to higher energy states of each clux sector contribute.  At low temperature (say $T_1$) only the low lying fermionic excitations over the  flux configurations plays the dominant role yielding a peak. However at higher temperature the fermionic excitations over each flux configurations are  being activated yielding another peak at temperature $T_2 > T_1$. 
When the non-Kitaev interaction are present, the fluxes no longer remains conserved. The direct implications of this is that the fermionic excitations can no linger be grouped together by flux quantum number and the low lying excitations over each flux sectors may not be separated by a gap and they may ultimately form a continuous band which can be continuously extended to the fermionic excitations at high energy. Thus we may lose the concept of two separate energy scales in the problem and find only a single peat in the specific heat. The presence of $J$ and $\Gamma$ destroys the above structure of the spectrum and fluxes are no longer conserved. This implies  any eigenstates can not be identified with a given flux quantum number. However one must note that a given energy eigenstate  can always be written in linear combinations of appropriate eigenstates of various flux quantum numbers. As the probability amplitude will in general be unequal for different flux quantum states, one can ask the question whether the notion of fermionic excitations and flux excitations still holds in the presence of non-Kitaev interactions. In this connections we have done extensive numerical analysis for all possible combinations of signs and magnitude of Kitaev interaction (K), Heisenberg interaction ($J$) and $\Gamma$ interactions and we find a great controllability of excitation spectra as revealed by specific heat as described below. \\
\indent
\begin{figure*}
	\includegraphics[width=1.2\columnwidth,height=!]{specific_heat_Gamma0_trimmed.pdf}
	\caption{Normalized specific heat $C_v$ in the presence of magnetic field is plotted. Each column denotes a particular cluster as mentioned at the top. At the right of each row, we mention the value of $(K,J)$ taken. In the inset $C_v$ for very large values of $J = \pm 2$ is shown. In the right the color legends show the value of magnetic field taken.}
\label{specificheatgammazero}
\end{figure*}
{\it{\small{General trends of specific heat :}}} In the Kitaev limit i.e when $K=\pm 1, J=\Gamma=0$, we find that specific heat has double peak structure for all the cluster except the four site cluster which has a weak shoulder like hump. Also the nature of specific heat is identical for $K=\pm 1$ for $J=\Gamma=0$  for a given cluster.  Presence of Heisenberg interaction acts differently depending on whether it is ferromagnetic or antiferromagnetic. For antiferromagnetic $J$ with $\gamma=0$, we find that only single peak in specific heat
exist but the position of the peaks can shift right or left depending on the magnitude of $J$. On the other hand for ferromagnetic $J$, after a critical $J_c$ we see that distinct appearances of two peak structure in specific heat. This mainly happens for $|J| > K$, though depending on the cluster size the two peak structure may appear or disappear as seen comparing the specific heat for six site and eight site cluster for $J=-0.5$ . On the other hand for ferromagnetic Kitaev interaction, the two peak structure is more frequent. For all ferromagnetic Heisenberg interaction it is present and for antiferromagnetic Heisenberg interaction the two peak structure occurs very occasionally only, for example for six site cluster with $J=1.0$. \\
\indent 
{\it{\small{Effect of $\Gamma$ on specific heat :}}} For $J=0$, we find that a ferromagnetic $\Gamma$ stabilizes the two peak structure  for $K=1$ and for $K=-1$, the effect of $\Gamma$(be it ferro or antiferro) destroys the two peak structure.  However when the magnitude of $J$ becomes larger than the magnitude of $\Gamma$, it is the Heisenberg interaction which controls the nature of specific heat. Interestingly at certain cases the effect of $\Gamma$ is seen to yield an additional third peak such as $(N,K,J,\Gamma)=(6,-1,1,0), (6,-1,1,0.5), (6,1,0,0.5), (12,1,-2,0)$  This implies the flux sectors further divided into two distinct energy scale in the systems. To elucidate the effect of $\Gamma$ we refer $(N,K,J,\Gamma)=(12,1,1.0,0.5)$ and $(12,-1,1,0.5)$ and we observe that in the former there is a distinct two peak structure which is absent on the later. Another universal feature is that for antiferromagnetic $J$ the peaks in the specific heat is more pronounced. Possibly the more distinct effect of $\Gamma$ can be seen for eight site cluster with $K=-1$. We see that for $\Gamma=0.5$, only $J=1$ yields a two peak structure. For $\Gamma=0$ all ferromagnetic $J$ (including $J=0$) has two peak structure but with $\Gamma=-0.5$ the two peak structure for negative $J$ becomes asymmetric and two peak structure reappear for positive $J=0.5, 1.0$. Thus we can say that $\Gamma$ can control the re-entrance of two peak structures in the specific heat.
\subsection{Effect of magnetic field on specific heat}
\indent
{\it \small{Ferromagnetic ($K=-1$) vs antiferromagnetic Kitaev ($K=1$) for $J=0$:}} The primary difference between
ferromagnetic and antiferromagnetic Kitaev interaction in the absence of Heisenberg interaction is that the two peaks in the case of antiferromagnetic Kitaev interaction are more frequent for many values of external magnetic field where for ferromagnetic Kitaev interaction the two peak structure only appears for small magnetic field, see Fig. \ref{specificheatgammazero} first row for antiferromagnetic Kitaev interaction and second row for ferromagnetic Kitaev interaction. This imply that the exact quantum eigenstates are renormalized for $K=1$ and $K=-1$ differently under the external magnetic field. The frequent occurrences of two peak structure implies that the eigenstates in the presence of external magnetic field are mildly modified in the case of antiferromagnetic Kitaev interaction. On the other hand for ferromagnetic Kitaev interactions, the effect of external magnetic field is more stronger to destroy the two peak structure implying confined phases. The broader wide peaks in the case of antiferromagnetic Kitaev interaction strengthen this understanding further that two peak structure mostly survives denoting robustness of deconfined phase in the presence of external magnetic field. With this  background understanding we now proceed to explain what happens when we take into consideration a finite antiferromagnetic Heisenberg interaction. \\
\indent
{\it \small{Ferromagnetic ($K=-1$) vs antiferromagnetic Kitaev ($K=1$) for $J=0.5$:}} In the case of antiferromagnetic Kitaev interaction the effect of having $J=0.5$ is that it makes the two peak structure
remain intact but with a  qualitative difference. If we compare Fig. \ref{specificheatgammazero} first row
(with $J=0$) and third row (with $J=0.5$), we observe that for $J=0.5$, the second peak is much more stronger
and the first peak is much weaker. Also the distance between the peaks has increased for $J=0.5$. This denotes
that antiferromagnetic Heisenberg interaction stabilizes the eigenspectrum determined by the conserved flux-sector qualitatively in the presence of magnetic field. Remarkably for ferromagnetic Kitaev interaction also we see re-entrance of two peak structure in the presence of antiferromagnetic Heisenberg interaction, see Fig. \ref{specificheatgammazero}, fourth row. We thus conclude that a antiferromagnetic Heisenberg interaction helps to stabilize the Kitaev-spin-liquid phase in comparison to absence of Heisenberg interaction in the presence of external magnetic field. \\
\indent
{\it \small{Ferromagnetic ($K=-1$) vs antiferromagnetic Kitaev ($K=1$) for $J=-0.5$:}} For the ferromagnetic Heisenberg interaction we observe that two peak structure are rarely occurring for antiferromagnetic Kitaev
interaction in comparison to $J=0.5$, compare Fig. \ref{specificheatgammazero} third row and fifth row.
Though the exact pattern of specific heat depends on the cluster sizes also. For example for all the cluster except twelve site, there are mostly a single peak. Similarly for the case of ferromagnetic Kitaev interaction we observe that there exists mostly single peak. This tells us that a ferromagnetic Heisenberg interaction works against the Kitaev-spin-liquid phase. It is interesting to note that an antiferromagnetic Heisenberg interaction
,in general, induces a  tendency to have a quantum paramagnet state which seems to favor a Kitaev-spin-liquid phase in comparison to a ferromagnetic Heisenberg interaction which induces a collinear order in the system which seems to destabilize the Kitaev-spin-liquid state. \\
\indent

 Having described the effect of various competing interactions of Kitaev-Heisenberg-$\Gamma$ model in finite temperature properties like magnetization,susceptibility and specific heat, we now cover a brief comparison with the existing results. The two peak structure of specific heat has been vindicated as the
hall mark signature of Kitaev spin liquid phase both theoretically and experimentally \cite{kataoka-jpsj-2020,andrade-2020,tathagata-2020,bachus-2021,balz-prb-2021,feng-2020,yogesh-2017,nasu-2015,yamaji-2016}.
The effect of magnetic field on these quantities have been very interesting some of which resonate with our
results. For example it has been shown that  the magnetic field help  to stabilize the 
magnetically ordered state before melting into a quantum paramagnetic state in $\alpha-RuCl_3$\cite{balz-prb-2021} where Kitaev interaction and Heisenberg interaction both has been considered ferromagnetic. In
accordance we find from the fourth row of Fig. \ref{specificheatgammazero} that when the $K$ and $J$ are both ferromagnetic only for zero or very small magnetic field we find two peak structure and an increased magnetic field stabilizes the one peak structure across all clusters. Similarly susceptibility shows qualitatively similar behavior found in Ref. \onlinecite{andrade-2020} though we have only shown the total susceptibilities. In fact in accordance with the previous results~\cite{andrade-2020}, we find in some cases, the susceptibility having maximum at some intermediate magnetic field and not at zero field, see Fig. \ref{susceptibility}, third, fourth and fifth row. However the combinations of $(J,K)$ is different than those taken in \onlinecite{andrade-2020} where $J$ and $K$ both are taken to be ferromagnetic and contain other complex interactions, it is interesting that the small cluster is able to capture the qualitative essence\cite{kataoka-jpsj-2020}. Similarly in another previous study~\cite{yamaji-2016} which considered a strong ferromagnetic Kitaev interaction and a week antiferromagnetic Kitaev interaction in the presence of next-nearest neighbor interaction, the specific heat is shown to change the width and height of the two peaks depending on the external magnetic field. This scenarios can be compared with Fig. \ref{specificheatgammazero}, fourth row, where across all the cluster similar behavior is observed. We comment that due to the limited cluster size, the  effect of next-nearest neighbor interaction is simulated easily.
\section{Discussion}
\label{discussion}
We now summarize our extensive investigation of some key feature of Kitaev-Hesineberg-$\Gamma$ model on finite size clusters. We started  with the inspection of eigenvalues of the system. It has been found that after a critical values of magnitude of Heisenberg interaction the spectrum varies linearly and conforms to the classical values for ferromagnetic Heisenberg interaction or appropriate singlet states for antiferromagnetic Heisenberg interaction. In the large ferromagnetic limit the energy density  varies as $E=-6 J + 2 K$ which does not depends on the cluster size. However, very interestingly, in the large anti-ferromagnetic Heisenberg interaction the asymptotic energy density depends on the cluster sizes which implies that the ground state is still a quantum ground state with large contribution of various possible  singlets which forms a linear combination due to Kitaev and $\Gamma$ interactions.  Near Kitaev limit, there are some branches of  eigenvalues which bend and do not change sign across zero as J is changed from negative to positive. Also there are large number of eigenenergies whose values vary linearly and  changes sign as the sign of $J$ is reversed.  Finally we have investigated how the eigenvalues vary with external magnetic field and we have seen remarkable closeness of the results found for large system. Though for ferromagnetic Kitaev interaction the effect of magnetic field is similar for all cluster, for antiferromagnetic $J$, the effect of magnetic field is shown to depend largely on the cluster sizes. This follows from the fact that the antiferromagnetic Heisenberg interaction induces  a significant quantum fluctuations which depend greatly on cluster sizes. \\
\indent
Next we investigated the bond-dependent correlation function on a given $z$ bond and looked at $<\sigma^z_i \sigma^z_j>$ and $<\sigma^y_i \sigma^y_j>$ correlations. In the Kitaev limit only $<\sigma^z_i \sigma^z_j>$ is non-zero on $z$-bond while $x-x$ and $y-y$ type correlation is zero.  In the presence of other non-Kitaev interaction $x-x$ and $y-y$ type correlation function becomes non-zero.  On the other hand in the Heisenberg limit, the bond-correlation function $<\sigma^{\alpha}_i \sigma^{\alpha}_j>$ is isotropic implying it is identical for $\alpha=x=y=z$. Thus while in the Kitaev limit the bond-dependent correlation function represents the extreme anisotropic character, the Heisenberg limit represents an isotropic character and we would expect a smooth transition between this as we increase the strength of Heisenberg  interaction. We observe that when $h_z$, the external magnetic field is absent this transition is indeed continuous. However in the presence of some intermediate values of external magnetic field the correlation function shows non-analytic and anisotropic character which manifests as presence of plateaus. This plateaus are only present for antiferromagnetic Heisenberg interaction implying that the different quantum ground state in the Kitaev limit and antiferromagnetic Heisenberg limit does not smoothly change from one another and there is discontinuity in it.\\
\indent
As mentioned earlier Kitaev model is known to contain two types of Majorana fermions named matter fermions and gauge fermions. In the Kitaev limit, two gauge fermions combine to constitute a conserved $Z_2$ flux operator. On the otherhand matter fermions constitute a non-interacting hopping problem coupled to this conserved $Z_2$ field. However in the presence of other non-Kitaev interaction  gauge fermions acquire dynamics. Various interesting consequences arise due to this such as confinement-deconfinement transition which is related to revival of spin-liquid phase in intermediate temperature. As we see that time evolution of Majorana fermions and gauge fermions follows oscillatory yet decaying pattern depending on the different interaction. There are various interesting pattern of this oscillations but must significant is their dependency on external magnetic field. External magnetic field is seen to greatly modulate the amplitude of oscillations and also the average magnitude of it. Depending on the strength of external magnetic field, it can freeze the gauge field without having any oscillations showing remarkable dependence of magnetic field. Interesting realization of quantum speed limit(QSL) phenomena also appear where probability of Majorana fermion vanishes from the initial value of unity. A detail analysis of this  QSL and their dependency on Heisenberg coupling and external magnetic field will be followed in future study.\\
\indent
Finally we have provided a detail analysis of finite temperature magnetization, susceptibility and specific heat for various parameters. It has been shown explicitly what is the salient difference between ferromagnetic  and antiferromagnetic Kitaev
interaction and what is effect of magnetic field and Heisenberg interaction (ferromagnetic and antiferromagnetic). Magnetization
shows clear distinction between ferro and antiferromagnetic Kitaev interaction. For low field and antiferromagnetic Kitaev interaction,
magnetization shows a dome like pattern which is absent for ferromagnetic Kitaev interaction. Also the zero temperature magnetization
have large dependency on external magnetic field. It can be evenly spaced from zero to a maximum or close to unity depending on whether Kitaev interaction is antiferromagnetic or antiferromagnetic. The Heisenberg interaction is shown to modulate the qualitative nature
of the magnetization further showing distinct quantum effects depending on whether it is antiferromagnetic or ferromagnetic. Similar
intriguing dependencies susceptibility and specific heat on various exchange coupling and external magnetic field has been observed. \\
\indent
With the recent intense effort to understand Kitaev-Heisenberg-$\Gamma$ system we think that our extensive analysis would serve as 
an useful reference. However we mention that all the aspect found here may not be readily generalized to thermodynamic system. Rather
it would be prudent to consider this as an interesting platform to quantify the quantum and thermal fluctuations pertaining to Kitaev-Heisenberg-$\Gamma$ system and use it for more better understanding of the thermodynamic system. For example in all the cases, universally, we have found that ferromagnetic and antiferromagnetic Kitaev interaction has distinct responses to external magnetic field
and Heisenberg interactions. It appears that in the case of antiferromagnetic Heisenberg interaction, the ground state has more quantum correlations irrespective of sign of Kitaev interaction which manifests itself in the plateaus like behavior in the correlation function, dome like shape in magnetization and stable two peak structure in specific heat. On the other hand a ferromagnetic Heisenberg
interaction is seen induce different characteristics. Thus a proper quantification of change of quantum fluctuation from the Kitaev limit to Heisenberg limit for both $J$ positive and negative would provide valuable insight to Kitaev physics. \\
\indent
 Kitaev-Heisenberg-$\Gamma$ system is shown to posses various collinear magnetic order at zero temperature which is destabilized at intermediate temperature towards a spin-liquid phase. A study of classical ground state on Kitaev cluster shown to exhibit interesting order by singularity mechanism within classical ground state space \cite{sarvesh-2020,anjana-2018}. An extension of this study to larger clusters and in the presence of Heisenberg interaction could provide valuable insight to explore new mechanism of quantum fluctuation~\cite{khatua-2021}. Further within the slave-fermion language implementing a PSG calculation\cite{wen-2002,yasir-2020,atanu-2023} for this finite size system utilizing various symmetries available and their response to various non-Kitaev interaction could be useful to compare the exact analysis done here and will be taken as future study. On a purely speculative way, we think that implementing a quench dynamics for the Kitaev system using external magnetic field or pressure could reveal interesting signatures of various phases obtained in Kitaev-Heisenberg-$\Gamma$ system.   


\section*{Acknowledgement} 
SMP acknowledges SAMKHYA (High-Performance Computing Facility provided by the Institute of Physics, Bhubaneswar) for the numerical computation. SM acknowledges support from the ICTP through the Associate's Programme (2020-2025). SM also thanks G. Baskaran and Nicola Seriani for interesting discussions.


\bibliography{bibfile}{}
\begin{thebibliography}{100}
\bibitem{kitaev-2006}A. Kitaev, Annals of Physics {\bf 321}, (2006) 2-111.
\bibitem{smandal-2007}G. Baskaran, S. Mandal, and R. Shankar, Phys. Rev. Lett. {\bf 98}, 247201 (2007).
\bibitem{kao-2021}Wen-Han Kao, Natalia B. Perkins, Annals of Physics {\bf 435} (2021) 168506.
\bibitem{nasu-2020} Joji Nasu and Yukitoshi Motome, Phys. Rev. B {\bf 102}, 054437 (2020). 
\bibitem{tathagata-2020}Tathagata Chowdhury, Achim Rosch, and Ralf Bulla, Phys. Rev. B {\bf 101}, 115133(2020).
\bibitem{tikhonov-2011} K. S. Tikhonov, M. V. Feigelman,  and A. Yu. Kitaev, Phys. Rev. Lett {\bf 106}, 067203(2011).
\bibitem{lunkin-2019}A. V. Lunkin, K. S. Tikhonov, M. V. Feigelman, Journal of Physics and Chemistry Solids, Volume 128, May 2019, Pages 130-137.
\bibitem{saptarshi-2011}Saptarshi Mandal, Subhro Bhattacharjee, K Sengupta, R Shankar, G Baskaran, Phys. Rev. B {\bf 84},155121 (2011).
\bibitem{animesh-2020}Animesh Nanda, Kusum Dhochak, and Subhro Bhattacharjee, Phys. Rev. B {\bf 102}, 235124 (2020).
\bibitem{animesh-2021}Animesh Nanda, Adhip Agarwala, and Subhro Bhattacharjee, Phys. Rev. B {\bf 104}, 195115 (2021).
\bibitem{knolle-2018}Johannes Knolle, Subhro Bhattacharjee, and Roderich Moessner, Phys. Rev. B {\bf 97}, 134432 (2018).
\bibitem{baskaran-2008}G. Baskaran, Diptiman Sen, and R. Shankar, Phys. Rev. B {\bf 78}, 115116 (2008).
\bibitem{samarth-2010}Samarth Chandra, Kabir Ramola, and Deepak Dhar, Phys. Rev E {\bf 82}, 031113(2010).
\bibitem{naveen-2008}Saptarshi Mandal and Naveen Surendran, Phys. Rev. B {\bf 79} 024426 (2008).
\bibitem{kivelson-2007}H. Yao and S. A. Kivelson, Phys. Rev. Lett. {\bf 99}, 247203 (2007).
\bibitem{vala-csl-2010}G. Kells, D. Mehta, J. K. Slingerland, and J. Vala, Phys. Rev. B {\bf 81}, 104429 (2010).
\bibitem{eschmann-2020} Tim Eschmann, Petr A. Mishchenko, Kevin O’Brien, Troels A. Bojesen, Yasuyuki Kato, Maria Hermanns, Yukitoshi Motome, and Simon Trebst, Phys. Rev. B {\bf 102}, 075125 (2020).
\bibitem{chalaupka-2010}Jiri Chaloupka, George Jackeli, and Giniyat Khaliullin Phys. Rev. Lett. {\bf 105}, 027204(2010).
\bibitem{simon-2022} Simon Trebst, Ciaran Hickey, Physics Reports {\bf 950} (2022)1-37.
\bibitem{takagi-2019} Hidenori Takagi, Tomohiro Takayama, George Jackeli, Giniyat Khaliullin \& Stephen E. Nagler, Nature Reviews Physics {\bf 1}, 264-280 (2019).
\bibitem{motome-2020}Yukitoshi Motome, Ryoya Sano, Seonghoon Jang, Yusuke Sugita1 and Yasuyuki Kato,  J. Phys.: Condens. Matter {\bf 32} 404001 (2020).
\bibitem{hermanns-2018}M. Hermanns, I. Kimchi, and J. Knolle, Annu. Rev. Condens. Matter Phys. {\bf 9}, 17–33 (2018).
\bibitem{abanerjee-2016}A. Banerjee, C. A. Bridges, J.-Q. Yan, A. A. Aczel, L. Li, M. B. Stone, G. E. Granroth,
M. D. Lumsden, Y. Yiu, J. Knolle, S. Bhattacharjee, D. L. Kovrizhin, R. Moessner, D. A. Tennant,
D. G. Mandrus, and S. E. Nagler, Nature Materials {\bf 15}, 733–740 (2016).
\bibitem{chern-2021}Li Ern Chern, Finn Lasse Buessen and Yong Baek Kim, npj Quantum Materials {\bf 6}, 33 (2016).
\bibitem{janssen-2017}Lukas Janssen, Eric C. Andrade,  and Matthias Vojta, Phys. Rev. B {\bf 96}, 064430 (2017).
\bibitem{czajka-2021} Peter Czajka, Tong Gao, Max Hirschberger, Paula Lampen-Kelley, Arnab Banerjee, Jiaqiang Yan, David G. Mandrus, Stephen E. Nagler and N. P. Ong, Nature Physics {\bf 17}, 915-919 (2021).
\bibitem{niravkumar-2019}Niravkumar D. Patel and Nandini Trivedi, PNAS {\bf 116} (25), 12199-12203 (2019).
\bibitem{haoxiang-nature-2021} Haoxiang Li, T. T. Zhang, A. Said, G. Fabbris, D. G. Mazzone, J. Q. Mandrus, Gabor B. Halasz, S. Okamoto, S. Murakami, M. P. M. Dean, H. N. Lee \& H Miao, Nature Communications {\bf 12}, 3513 (2012).
\bibitem{balz-prb-2021}C. Balz, L. Janssen, P. Lampen-Kelley, A. Banerjee, Y. H. Liu, J.-Q. Yan, D. G. Mandrus, M. Vojta, and S. E. Nagler, Phys. Rev. B {\bf 103}, 174417 (2021).
\bibitem{berke-prb-2020}Christoph Berke, Simon Trebst, and Ciarán Hickey, Phys. Rev. B {\bf 101}, 214442 (2020).
\bibitem{plessis-2020} D. Reig-i-Plessis, T. A. Johnson, K. Lu, Q. Chen, J. P. C. Ruff, M. H. Upton, T. J. Williams, S. Calder, H. D. Zhou, J. P. Clancy, A. A. Aczel, and G. J. MacDougall, Phys. Rev. Mat. {\bf 4}, 124407 (2020).
\bibitem{wulferding-2020}Dirk Wulferding, Youngsu Choi, Seung-Hwan Do, Chan Hyeon Lee, Peter Lemmens, Clement Faugeras, Yann Gallais \& Kwang-Yong Choi,
Nature Communications {\bf 11} 1603 (2020).
\bibitem{sananda-2019}Sananda Biswas, Ying Li, Stephen M. Winter, Johannes Knolle, and Roser Valenti, Phys. Rev. Lett {\bf 123}, 237201 (2019).
\bibitem{sananda-2021}V. Leeb, K. Polyudov, S. Mashhadi, S. Biswas,  Roser Valenti,  M. Burghard, and J. Knolle, Phys. Rev. Lett, {\bf 126}, 097201 (2021).
\bibitem{anderson}P. W. Anderson, Science, {\bf 177}, 393 (1972).
\bibitem{fifty-years}Steven Strogatz, Sara Walker, Julia M. yeomans, Corina tarnita, Elsa Arcaute, Manlio De Domenico, Oriol Artime and Kwang-II Goh, Nature Reviews Physics {\bf 4}, 508-510 (2022). 
\bibitem{savary-review}Lucile Savary and Leon Balents, Rep. Prog. Phys. {\bf 80} 016502 (2017).
\bibitem{kataoka-jpsj-2020}Kota Kataoka, Daigoro Hirai, Takeshi Yajima, Daisuke Nishio-Hamane, Rieko Ishii, Kwang-Yong Choi, Dirk Wulferding, Peter Lemmens, Shunichiro Kittaka,Toshiro Sakakibara, Hajime Ishikawa, Akira Matsuo, Koichi Kindo, and Zenji Hiroi,  J. Phys. Soc. Jpn. {\bf 89}, 114709 (2020).
\bibitem{wang-2020} Yiping Wang, Gavin B. Osterhoudt, Yao Tian, Paige Lampen-Kelley, Arnab Banerjee, Thomas Goldstein, Jun Yan,
Johannes Knolle, Huiwen Ji, Robert J. Cava, Joji Nasu, Yukitoshi Motome, Stephen E. Nagler, David Mandrus \& Kenneth S. Burch, npj Quantum Materials {\bf 5}, 14 (2020).
\bibitem{wang-2021} Shi Wang, Zhongyuan Qi, Bin Xi, Wei Wang, Shun-Li Yu, and Jian-Xin Li, Phys. Rev. B {\bf 103}, 054410 (2021).
% \cite{berke-prb-2020,nasu-motome-2015,koga-nasu-2019,li-2021-ncom }
\bibitem{nasu-motome-2015} Joji Nasu, Masafumi Udagawa, and Yukitoshi Motome, Phys. Rev. B {\bf 92}, 115122 (2015).
\bibitem{koga-nasu-2019} Akihisa Koga and Joji Nasu, Phys. Rev. B {\bf 100}, 100404(R) (2019).
\bibitem{li-2021-ncom}Han Li, Hao-Kai Zhang, Jiucai Wang, Han-Qing Wu, Yuan Gao, Dai-Wei Qu,
Zheng-Xin Liu, Shou-Shu Gong $\&$ Wei Li, Nature Communications v{\bf 12},  4007 (2021).
\bibitem{zhu-2018} Zheng Zhu, Itamar Kimchi, D. N. Sheng and Liang Fu, Phys. Rev. B {\bf 97}, 241110(R) (2018).
\bibitem{holzmann-2016}Markus Holzmann, Raymond C. Clay III, Miguel A. Morales, Norm M. Tubman, David M. Ceperley, and Carlo Pierleoni, Phys. Rev B {\bf 94}, 035126 (2016).
\bibitem{han-li-2023} Han Li, Wei Li and Gang Su, Phys. Rev. B {\bf 107}, 115124 (2023).
\bibitem{yilmaz-2022}F. Yilmaz, A. P. Kampf, and S. K. Yip,  Phys. Rev. Research {\bf 4}, 043024 (2022).
\bibitem{khwang-2022} Kyusung Hwang, Ara Go, Ji Heon Seong, Takasada Shibauchi and Eun-Gook Moon, Nature Communications {\bf 13}, 323 (2022). 
\bibitem{yamaji-2016} Y. Yamaji, T. Suzuki, T. Yamada, S. I. Suga, N. Kawashima, and M. Imada, Phys. Rev. B {\bf 93}, 174425 (2016).
\bibitem{feng-2022}Shi Feng, Adhip Agarwala, Subhro Bhattacharjee, Nandini Trivedi, condmat, arXiv:2206.12990v2
\bibitem{shuyi-2023}Shuyi Li, Han Yan, and Andriy H. Nevidomskyy, condmat arXiv:2301.07401v2
\bibitem{smit-2020}R. L. Smit, S. Keupert, O. Tsyplyatyev, P. A. Maksimov, A. L. Chernyshev, and P. Kopietz, Phys. Rev. B {\bf 101}, 054424 (2020).
\bibitem{kimchi-2014}Itamar Kimchi, James G. Analytis, and Ashvin Vishwanath, Phys. Rev. B {\bf 90}, 205126 (2014).
\bibitem{young-baek-kim-2015} Eric Kin-Ho Lee and Yong Baek Kim, Phys. Rev. B {\bf 91}, 064407(2015).
\bibitem{song-2016} Xue-Yang Song, Yi-Zhuang You, and Leon Balents, Phys. Rev. Lett {\bf 117}, 037209 (2016).
\bibitem{yang-2020}Fan Yang, Kirill Plekhanov, and Karyn Le Hur, Phys. Rev. Research {\bf 2}, 013005(2020). 
\bibitem{lieb-1994}E. H. Lieb, Phys. Rev. Lett. {\bf 73}, 2158 (1994).
\bibitem{nandini-2019}David C. Ronquillo, Adu Vengal and Nandini Trivedi, Phys. Rev. B {\bf 99}, 140413(R) (2019).
\bibitem{liang-2018}Shuang Liang, Bo-Sen He, Zhao-Yang Dong, Wei Chen, Jian-Xin Li, and Qiang-Hua Wang, Phys. Rev. B {\bf 98}, 104410 (2018).
\bibitem{mandaljpa}Saptarshi Mandal, R Shankar, G Baskaran, Journal of Physics A: Math. and Theor. {\bf 45}, 335304 (2012).
\bibitem{nasu-motome-2019}Joji Nasu and Yukitoshi Motome, Phys. Rev. Research  {\bf 1}, 033007 (2019).
\bibitem{yamada-2020} Masahiko G. Yamada, npj Quantum Materials {\bf 5}, 82 (2020).
\bibitem{deffner-2017}Sebastian Deffner and Steve Campbell, J. Phys. A: Math. Theor. {\bf 50}, 453001 (2017).
\bibitem{gal-ness-2022}Gal Ness, Andrea Alberti and Yoav Sagi, Phys. Rev. Lett. {\bf 129}, 140403 (2022).
\bibitem{pati-2022}Brij Mohan  and Arun Kumar Pati, Phys. Rev. A {\bf 106}, 042436 (2022).
\bibitem{arindam-2022}Srishty Aggarwal, Subhashish Banerjee, Arindam Ghosh and Banibrata Mukhopadhyay,  New J. Phys. {\bf 24}, 085001 (2022).
\bibitem{bachus-2021}S. Bachus, D. A. S. Kaib, A. Jesche, V. Tsurkan, A. Loidl,  S. M. Winter,
A. A. Tsirlin, R. Valenti, and P. Gegenwart, Phys. Rev. B {\bf 103}, 054440 (2021).
\bibitem{andrade-2020} Eric C. Andrade, Lukas Janssen, and Matthias Vojta, Phys. Rev. B {\bf 102}, 115160 (2020).
\bibitem{feng-2020}Kexin Feng, Natalia B. Perkins , and F. J. Burnell, Phys. Rev. B {\bf 102}, 224402 (2020).
\bibitem{yogesh-2017}Kavita Mehlawat, A. Thamizhavel, and Yogesh Singh, Phys. Rev. B {\bf 95}, 144406 (2017).
\bibitem{nasu-2015}J. Nasu, M. Udagawa and Y. Motome, Phys. Rev. B {\bf 92}, 115122(2015).
\bibitem{anjana-2018} Anjana M. Samarakoon, Gideon Wachtel, Youhei Yamaji, D. A. Tennant, Cristian D. Batista, and Yong Baek Kim, Phys. Rev. B {\bf 98}, 045121 (2018).
\bibitem{sarvesh-2020} Sarvesh Srinivasan, Subhankar Khatua,G. Baskaran, and R. Ganesh, Phys. Rev. Research {\bf 2}, 023212(2020).
\bibitem{khatua-2021}Subhankar Khatua, Sarvesh Srinivasan, and R. Ganesh, Phys. Rev. B {\bf 103}, 174412 (2021).
%wen-2002,yasir-2020,atanu-2023
\bibitem{wen-2002}X.-G. Wen, Phys. Rev. B {\bf 65}, 165113 (2002).
\bibitem{yasir-2020}J. Sonnenschein, A. Chauhan, Y. Iqbal, and J. Reuther, Phys. Rev. B {\bf 102}, 125140 (2020).
\bibitem{atanu-2023}Atanu Maity, Francesco Ferrari, Ronny Thomale, Saptarshi Mandal, and Yasir Iqbal, Phys. Rev. B {\bf 107}, 134438 (2023).
%finite-size-system-study references upto above done
%\bibitem{yamada-2021}Masahiko G. Yamada, and Satoshi Fujimoto, Phys. Rev. Lett {\bf 127}, 047201 (2021)
%\bibitem{natori-2020}Willian M. H. Natori and Johannes Knolle, Phys. Rev. Lett {\bf 125}, 067201 (2020)
%\bibitem{takahasi-2021}Masahiro O. Takahashi, Masahiko G. Yamada, Daichi Takikawa, Takeshi Mizushima, and Satoshi Fujimoto, Phys. Rev. Research {\bf 3}, 023189 (2021)
%\bibitem{zhang-2021}Emily Z. Zhang, Li Ern Chern, and Yong Baek Kim Phys. Rev. B {\bf 103}, 174402
%%Experiment
%\bibitem{young-baek-kim-2015} Eric Kin-Ho Lee and Yong Baek Kim, Phys. Rev. B {\bf 91}, 064407(2015).
%\bibitem{klocke-2021}Kai Klocke, David Aasen, Roger S. K. Mong, Eugene A. Demler, and Jason Alicea, Phys. Rev. Lett {\bf 126}, 177204(2021)
%\bibitem{dantas-2022} Vitor Dantas and Eric C. Andrade, Phys. Rev. Lett {\bf 129}, 037204 (2022).
%\bibitem{fraser-1996}Louisa M. Fraser and W. M. C. Foulkes, G. Rajagopal, R. J. Needs, S. D. Kenny, and A. J. Williamson, Phys. Rev. B {\bf 53},1814 (1996).
%\bibitem{tomoya-2022}Tomoya Iwashita, Masaaki Nagao, Akira Yoshimori, Masahide Terazima, Ryo Akiyama, Chem. Phys. Lett. {\bf 807}, 140096 (2022).
%yogesh-2017,nasu-2015,yamaji-2016
%quantum-speed-limit  deffner-2017,gal-ness-2022,pati-2022,arindam-2022
%khwang-2022,feng-2022,han-li-2023,yilmaz-2022,shuyi-2023,smit-2020,kimchi-2014,yamaji-2016,young-baek-kim-2015,wang-2021
%shubhro animesh-2020, animesh-2021, knolle-2018
\end{thebibliography}
\section{Appendix}
\label{appendix}
\appendix
\newcounter{defcounter}
\setcounter{defcounter}{0}

\vspace{20pt}
\section{Eigenstates of Kitaev-Heisenberg 4-site cluster}{\label{App:A}}

In 4-site cluster, with only Kitaev interaction the Hamiltonian is easily diagonalized and the eigenvalues are $ \{\pm 2K_x, \pm 2K_y, \pm 2K_z, \pm 2\sqrt{K_x^2+K_y^2+K_z^2}\}$ with each energy level having two-fold degeneracy. The Hamiltonian can be diagonalized into two sectors one with even number of spin up and the other one with odd number of up. Two degenerate ground states in Kitaev limit, are:
\begin{eqnarray}
	\ket{{\rm GS}_1}&=&K_y\left( \ket{\uparrow\uparrow\uparrow\uparrow}+\ket{\downarrow\downarrow\downarrow\downarrow}\right)-K_x\left(\ket{\uparrow\downarrow\uparrow\downarrow}+\ket{\downarrow\uparrow\downarrow\uparrow}\right)\nonumber\\
	&&+\left( K_z+\tilde{K}\right) \left( \ket{\uparrow\downarrow\downarrow\uparrow}+\ket{\downarrow\uparrow\uparrow\downarrow} \right), \nonumber\\
	\ket{{\rm GS}_2}&=&K_y(\ket{\uparrow\uparrow\downarrow\downarrow}+\ket{\downarrow\downarrow\uparrow\uparrow})+K_x(\ket{\uparrow\downarrow\downarrow\uparrow}+\ket{\downarrow\uparrow\uparrow\downarrow})\nonumber\\
	&&+(K_z-\tilde{K})(\ket{\uparrow\downarrow\uparrow\downarrow}+\ket{\downarrow\uparrow\downarrow\uparrow}),\nonumber\\
	{\rm where}&&\tilde{K}=\sqrt{K_x^2+K_y^2+K_z^2}.
\end{eqnarray}
The spin-spin correlations (CR) can also be written in closed analytical form. The non vanishing correlations are:
\begin{eqnarray}
	{\rm CR}^{x-x}&=&-4\frac{K_x}{K_y^2}(K_z+\sqrt{K_x^2+K_y^2+K_z^2})/N_0;\nonumber\\
	{\rm CR}^{y-y}&=&-4\frac{1}{K_y}(K_z+\sqrt{K_x^2+K_y^2+K_z^2})/N_0;\nonumber\\
	{\rm CR}^{z-z}&=&-4\frac{K_z}{K_y^2}(K_z+\sqrt{K_x^2+K_y^2+K_z^2})/N_0;\nonumber\\
	N_0&=&\sqrt{ 2+2\frac{K_x^2}{K_y^2}+\frac{2}{K_y^2}(K_z+\sqrt{K_x^2+K_y^2+K_z^2})^2}.\nonumber\\
\end{eqnarray}
 With the application of Heisenberg interaction $J$ the Hamiltonian looses some of the symmetries and the solution is more complex. The eigenvalues of the odd-particle sector are $-2J\pm2K_y,-2J-2K_y,-2J\pm2K_x,-2J-2K_x,6J+2K_y,6J+2K_x$ and the even-particle sector has eigenvalues $6J+2K_z,-2J\pm 2K_z,-2J-2K_z,\pm2\left(6J^2+\sum_{\alpha}^{}(J+K_\alpha)^2\pm J\sqrt{8\sum_{\alpha}^{}(K_\beta-K_\gamma)^2}\right)^{1/2}$, where each of $\alpha,\beta,\gamma$ can take values $x,y,z$ and $\alpha\neq\beta\neq\gamma$.\\
\\
In the presence of only Heisenberg interaction $J$ (i.e., $K$=0), the cluster is SU(2) symmetric.  In this limit, the ground state is a singlet state of the form $\frac{1}{2}\left(\ket{\uparrow_1\downarrow_2}-\ket{\downarrow_1\uparrow_2}\right)\times\left(\ket{\uparrow_3\downarrow_4}-\ket{\downarrow_3\uparrow_4}\right)$. When we introduce external magnetic field $h_z$, we find presence of $\ket{\uparrow\uparrow\downarrow\downarrow},\ket{\downarrow\downarrow\uparrow\uparrow}$ states in the ground state with some amplitude. When we increase the magnitude of $h_z$, this amplitude changes but only these six states make the ground state.  For ferromagnetic $J$, even a small magnetic field yields a large contribution of  $\ket{\uparrow\uparrow\uparrow\uparrow}$ to the ground state. The 4-site cluster in the presence of Kitaev and Heisenberg interaction, can be solved analytically. The Hamiltonian becomes block-diagonalized of two blocks for even and odd particle sector. The basis states used for these two sectors are:
\begin{eqnarray}
	\ket{\rm even}=
	\begin{pmatrix}
		\ket{\uparrow\uparrow\uparrow\uparrow}\\
		\ket{\uparrow\downarrow\downarrow\uparrow}	\\
		\ket{\downarrow\uparrow\uparrow\downarrow}	\\
		\ket{\uparrow\uparrow\downarrow\downarrow}	\\
		\ket{\downarrow\downarrow\uparrow\uparrow}	\\
		\ket{\uparrow\downarrow\uparrow\downarrow}	\\
		\ket{\downarrow\uparrow\downarrow\uparrow}	\\
		\ket{\downarrow\downarrow\downarrow\downarrow}
	\end{pmatrix};
	&&
	\ket{\rm odd}=
	\begin{pmatrix}
		\ket{\downarrow\uparrow\uparrow\uparrow}\\
		\ket{\uparrow\downarrow\uparrow\uparrow}\\
		\ket{\uparrow\uparrow\downarrow\uparrow}\\
		\ket{\uparrow\uparrow\uparrow\downarrow}\\
		\ket{\uparrow\downarrow\downarrow\downarrow}\\
		\ket{\downarrow\uparrow\downarrow\downarrow}\\
		\ket{\downarrow\downarrow\uparrow\downarrow}\\
		\ket{\downarrow\downarrow\downarrow\uparrow}
	\end{pmatrix};
\end{eqnarray}
and the odd particle sector is diagonalized in the rotated basis:

{\small
\begin{eqnarray}
	\frac{1}{\sqrt{8}}
\begin{pmatrix}
	1 & -1 & -1 & 1 & -1 & 1 & 1 & -1 \\
	1 & 1 & 1 & 1 & -1 & -1 & -1 & -1 \\
	0 & \sqrt{2} & -\sqrt{2} & 0 & 0 & -\sqrt{2} & \sqrt{2} & 0 \\
	\sqrt{2} & 0 & 0 & -\sqrt{2} & -\sqrt{2} & 0 & 0 & \sqrt{2} \\
	-\sqrt{2} & \sqrt{2} & 0 & 0 & -\sqrt{2} & \sqrt{2} & 0 & 0 \\
	1 & 1 & -1 & -1 & 1 & 1 & -1 & -1 \\
	1 & 1 & 1 & 1 & 1 & 1 & 1 & 1 \\
	0 & 0 & \sqrt{2} & -\sqrt{2} & 0 & 0 & \sqrt{2} & -\sqrt{2} \\
\end{pmatrix}
\begin{pmatrix}
	\ket{\downarrow\uparrow\uparrow\uparrow}\\
	\ket{\uparrow\downarrow\uparrow\uparrow}\\
	\ket{\uparrow\uparrow\downarrow\uparrow}\\
	\ket{\uparrow\uparrow\uparrow\downarrow}\\
	\ket{\uparrow\downarrow\downarrow\downarrow}\\
	\ket{\downarrow\uparrow\downarrow\downarrow}\\
	\ket{\downarrow\downarrow\uparrow\downarrow}\\
	\ket{\downarrow\downarrow\downarrow\uparrow}	
\end{pmatrix}.
\end{eqnarray}
}
The corresponding eigenvalues are $-2J+2K_y,6J+2K_y,-2J-2K_y,-2J-2K_y,-2J-2K_x,-2J+2K_x,6J+2K_x,-2J-2K_x$. When the even-sector basis is rotated to a new one:
\begin{eqnarray}
	\begin{pmatrix}
	\ket{\uparrow\uparrow\uparrow\uparrow}\\
	\ket{\uparrow\downarrow\downarrow\uparrow}	\\
	\ket{\downarrow\uparrow\uparrow\downarrow}	\\
	\ket{\uparrow\uparrow\downarrow\downarrow}	\\
	\ket{\downarrow\downarrow\uparrow\uparrow}	\\
	\ket{\uparrow\downarrow\uparrow\downarrow}	\\
	\ket{\downarrow\uparrow\downarrow\uparrow}	\\
	\ket{\downarrow\downarrow\downarrow\downarrow}
\end{pmatrix}
&
\rightarrow
&
\frac{1}{\sqrt{2}}
\begin{pmatrix}
(1+\Delta)\ket{\uparrow\uparrow\uparrow\uparrow}\\
(1+\Delta)\ket{\uparrow\downarrow\downarrow\uparrow}\\
(1+\Delta)\ket{\uparrow\uparrow\downarrow\downarrow}\\
(1+\Delta)\ket{\uparrow\downarrow\uparrow\downarrow}\\
(1-\Delta)\ket{\uparrow\uparrow\uparrow\uparrow}\\
(1-\Delta)\ket{\uparrow\downarrow\downarrow\uparrow}\\
(1-\Delta)\ket{\uparrow\uparrow\downarrow\downarrow}\\
(1-\Delta)\ket{\uparrow\downarrow\uparrow\downarrow}\\
\end{pmatrix},
\end{eqnarray}
where $\Delta=\sigma_{1}^{x}\sigma_{2}^{x}\sigma_{3}^{x}\sigma_{4}^{x}$, the even-sector of the Hamiltonian is block diagonalized in two blocks. The lower four are eigenstates (which are part of the second of the two blocks) with eigenvalues $6J+2K_z,-2J-2K_z,-2J-2K_z,-2J+2K_z$. The upper four basis states reduce the Hamiltonian to a $4\times4$ block:
\begin{equation}
	\begin{pmatrix}
		6J+2K_z&-2K_y&2K_x&0\\
		-2K_y&-2J-2K_z&4J&4J+2K_x\\
		2K_x&4J&-2J-2K_z&4J+2K_y\\
		0&4J+2K_x&4J+2K_y&-2J+2K_z
	\end{pmatrix},
\end{equation}
and this gets diagonalized in the basis of,
{
\small
\begin{eqnarray}
	\frac{(1+\Delta)}{\sqrt{2}}\left(\upsilon_1^{(1)} \ket{\uparrow\uparrow\uparrow\uparrow}+\upsilon_1^{(2)}\ket{\uparrow\downarrow\downarrow\uparrow}+\upsilon_1^{(3)}\ket{\uparrow\uparrow\downarrow\downarrow}+\upsilon_1^{(4)}\ket{\uparrow\downarrow\uparrow\downarrow}\right)\nonumber \\
	\frac{(1+\Delta)}{\sqrt{2}}\left(\upsilon_2^{(1)} \ket{\uparrow\uparrow\uparrow\uparrow}+\upsilon_2^{(2)}\ket{\uparrow\downarrow\downarrow\uparrow}+\upsilon_2^{(3)}\ket{\uparrow\uparrow\downarrow\downarrow}+\upsilon_2^{(4)}\ket{\uparrow\downarrow\uparrow\downarrow}\right) \nonumber\\
	\frac{(1+\Delta)}{\sqrt{2}}\left(\upsilon_3^{(1)} \ket{\uparrow\uparrow\uparrow\uparrow}+\upsilon_3^{(2)}\ket{\uparrow\downarrow\downarrow\uparrow}+\upsilon_3^{(3)}\ket{\uparrow\uparrow\downarrow\downarrow}+\upsilon_3^{(4)}\ket{\uparrow\downarrow\uparrow\downarrow}\right)\nonumber \\
	\frac{(1+\Delta)}{\sqrt{2}}\left(\upsilon_4^{(1)} \ket{\uparrow\uparrow\uparrow\uparrow}+\upsilon_4^{(2)}\ket{\uparrow\downarrow\downarrow\uparrow}+\upsilon_4^{(3)}\ket{\uparrow\uparrow\downarrow\downarrow}+\upsilon_4^{(4)}\ket{\uparrow\downarrow\uparrow\downarrow}\right),\nonumber\\
\end{eqnarray}
}
\begin{widetext}
		\begin{eqnarray}
		\upsilon_1&=&\bot_1
		\begin{pmatrix}
			\left(-3J^2(K_{\Sigma}-3K_z)-2J K_{R} (K_z-\lambda_{-})-J\left(-K_{R}^2-K_z^2+K_x K_y+6J K_{R}\right)+J
			\lambda_{-} (K_{\Sigma}-3K_z)\right) \\
			3J^2(K_x-K_y)+2J K_y K_{R}+J\left(K_{R}^2-K_z^2+K_x K_y\right)-J \lambda_{-}(K_x-K_y)
			\\
			3J^2(K_x-K_y)-2J K_x K_{R} - J\left(K_{R}^2-K_z^2+K_x K_y\right)-J \lambda_{-}(K_x-K_y)
			\\
			J(K_x-K_y)(3J+K_{\Sigma}-\lambda_{-})
		\end{pmatrix}\\
		\upsilon_2&=&\bot_2
		\begin{pmatrix}
			\left(-3J^2(K_{\Sigma}-3K_z)-2J K_{R} (K_z+\lambda_{-})-J\left(-K_{R}^2-K_z^2+K_x K_y+6J K_{R}\right)-J
			\lambda_{-} (K_{\Sigma}-3K_z)\right) \\
			3J^2(K_x-K_y)+2J K_y K_{R}+J\left(K_{R}^2-K_z^2+K_x K_y\right)+J \lambda_{-}(K_x-K_y)
			\\
			3J^2(K_x-K_y)-2J K_x K_{R} - J\left(K_{R}^2-K_z^2+K_x K_y\right)+J \lambda_{-}(K_x-K_y)
			\\
			J(K_x-K_y)(3J+K_{\Sigma}+\lambda_{-}) \\
		\end{pmatrix}\\
		\upsilon_3&=&\bot_3
		\begin{pmatrix}
			\left(-3J^2(K_{\Sigma}-3K_z)+2J K_{R} (K_z-\lambda_{+})+J\left(K_{R}^2+K_z^2-K_x K_y+6J K_{R}\right)+J
			\lambda_{+} (K_{\Sigma}-3K_z)\right) \\
			3J^2(K_x-K_y)-2J K_y K_{R}+J\left(K_{R}^2-K_z^2+K_x K_y\right)-J \lambda_{+}(K_x-K_y)
			\\
			3J^2(K_x-K_y)+2J K_x K_{R} - J\left(K_{R}^2-K_z^2+K_x K_y\right)-J \lambda_{+}(K_x-K_y)
			\\
			J(K_x-K_y)(3J+K_{\Sigma}-\lambda_{+}) \\
		\end{pmatrix}\\
		\upsilon_4&=&\bot_4
		\begin{pmatrix}
			\left(-3J^2(K_{\Sigma}-3K_z)+2J K_{R} (K_z+\lambda_{+})+J\left(K_{R}^2+K_z^2-K_x K_y+6J K_{R}\right)-J
			\lambda_{+} (K_{\Sigma}-3K_z)\right) \\
			3J^2(K_x-K_y)-2J K_y K_{R}+J\left(K_{R}^2-K_z^2+K_x K_y\right)+J \lambda_{+}(K_x-K_y)
			\\
			3J^2(K_x-K_y)+2J K_x K_{R} - J\left(K_{R}^2-K_z^2+K_x K_y\right)+J \lambda_{+}(K_x-K_y)
			\\
			J(K_x-K_y)(3J+K_{\Sigma}+\lambda_{+}) \\	
		\end{pmatrix}
	\end{eqnarray}
\end{widetext}	
with proper normalization $\bot_1,\bot_2,\bot_3,\bot_4$. $\upsilon_i^{(j)}$ is the $j^{\rm th}$ component of $\upsilon_i$. Here we have used the following abbreviations:
\begin{eqnarray}
K_{\Sigma}&=&K_x+K_y+K_z;\\
\tilde{K}&=&\sqrt{K_x^2+K_y^2+K_z^2};\\
K_R&=&\sqrt{\tilde{K}^2-K_xK_y-K_yK_z-K_zK_x};\\
\lambda_{\pm}&=&\sqrt{9J^2+\tilde{K}^2+2JK_{\Sigma}\pm4JK_R}.
\end{eqnarray}
In this basis, the $4\times4$ block is diagonalized with eigenvalues $-2\lambda_{-},2\lambda_{-},-2\lambda_{+},2\lambda_{+}$.


\begin{comment}
\end{comment}
\end{document}
